\documentclass{article}
\usepackage{amsmath}
\usepackage{amsfonts}
\usepackage{bbold}
\usepackage{xcolor}
\usepackage{hyperref}
\usepackage[italian]{babel}

\definecolor{blu}{HTML}{ddeaff}
\definecolor{orange}{HTML}{ffead0}


\title{
  \textbf{Quesiti svolti di Laboratorio3} \\
  \large Universit\'a di Pisa 2025-2026
}
\author{your dear green}

\begin{document}
\maketitle
\begin{abstract}
  Compilazione di prove di esame del corso di \textit{Laboratorio 3}
  per la laurea triennale in Fisica all'Universit\'a di Pisa tenuto da
  Donato Niccol\'o e Chiara Maria Angela Roda.
\end{abstract}

% TODO
% 1) Presentare le funzioni di trasferimento in FORMA CANONICA.

\section*{Ringraziamenti}
\paragraph{Camilla Maccari} mi sono appoggiato molto sui
suoi esercizi svolti nei vari tentativi di superare l'esame. I
seguenti sono poco pi\'u di una rifinitura e qualche completamento
delle note di Camilla.

\paragraph{Vittorio Scirli} lo ringrazio molto per avermi fatto scopire
il software \textit{falstad}, che ha permesso di simulare i circuiti con
facilit\'a e di dare un criterio di compatibilit\'a tra i risultati
da me ottenuti e la ``realt\'a'' (simulata da falstad). Lo ringrazio
anche per avermi aiutato in generale con questa materia.

\paragraph{Marco Nencioni} lo ringrazio per essere stato il mio
compagno di laboratorio dal \textit{primo} fino al \textit{terzo}
anno. Siamo passati, da non saper fare il fit di una retta, a non capire
la differenza fra serie e parallelo fino a non capire come fare un
semaforo.

\section*{Leggimi}
I seguenti esercizi non sono guidati dal lato di vista teorico, il
loro scopo \'e quello di dare un riferimento durante la preparazione
dell'esame alla luce della teoria \textit{gi\'a imparata} da materiale
esterno a questo; tuttavia pu\'o essere anche utile questo testo come
fonte di esempi per accompagnare l'apprendimento.
Detto questo buona fortuna!


\pagebreak

\section*{Bandpass VCVS}

\noindent\colorbox{orange}{
  \begin{minipage}{1.0\linewidth}
    \paragraph{Richiesta (a)}
    Dimostrare, sulla base di semplici argomenti di natura qualitativa, che il circuito proposto
    svolga effettivamente la funzione di filtro passa-banda, discutendone la risposta ai limiti
    del dominio delle frequenze reali.
  \end{minipage}
}
\paragraph{Risposta (a)}
La tensione d'ingresso $V_S$ attraversa un arrangiamento di partitori configurati come
un filtro passa basso (passivo) e un filtro passa alto (passivo) collegati a cascata;
il filtro passa basso ha un ramo in comune con la retroazione dell'opamp.
L'opamp segue lo schema non invertente.

Al limite del dominio delle frequenze reali si manifestano le seguenti condizioni
\begin{itemize}
\item
  Per frequenze $f \to 0$ i condensatori si comportanto come rami aperti il nodo B
  \'e messo a terra e nel regime di operazione lineare
  \begin{equation*}
    V_{OUT} = A_d(0-V_-) = A_d(0-kV_{OUT})
  \end{equation*}
  la tensione all'ingresso dell'ingresso invertente ha un fattore di partizione positivo
  $k$ (non \'e necessario calcolarlo)
  \begin{equation*}
    A_d(0-kV_{OUT}) \implies V_{OUT} = \frac{1}{1+kA_d}
  \end{equation*}
  per guadagni differenziali elevati la tensione $\colorbox{blu}{\boxed{V_{OUT} \to 0}}$
\item
  Per frequenze $f \to \infty$ i condensatori si comportanto come corto-circuiti il nodo A
  \'e messo a terra, la tensione al nodo B \'e nulla quindi, per un argomento analogo al caso
  $f\to 0$, per guadagni differenziali elevati la tensione $\colorbox{blu}{\boxed{V_{OUT} \to 0}}$.
\end{itemize}

\noindent\colorbox{orange}{
  \begin{minipage}{1.0\linewidth}
    \paragraph{Richiesta (b)}
    Derivare analiticamente la funzione di trasferimento complessa Av(s)=VOUT(s)/VS(s)
    nell’ipotesi che l’operazionale sia ideale (\textit{suggerimento: si determini la tensione nel
      punto A del circuito in termini di VS e VOUT; indipendentemente, si esprimano VB e quindi
      VA in termini di Vout assumendo che l'opAmp funzioni in regime lineare ed utilizzando il
      principio di cortocircuito virtuale; si eguaglino infine le espressioni ottenute per ottenere
      la funzione di trasferimento}). Per semplicit\'a di calcolo, si assuma inoltre (ipotesi
    compatibile con i valori e le incertezze nominali dei componenti) che R1 = R3 = R, R2 =
    2 R, C1 = C2 = C, $R_B$ = 5/3 $R_A$.
  \end{minipage}
}
\paragraph{Risposta (b)}
La richiesta porta con se un suggerimento, seguendolo per il nodo A vale
\begin{equation*}
  \frac{V_S - V_A}{R} = sCV_A + sC(V_A - V_B) + \frac{V_A - V_{OUT}}{R}
\end{equation*}
quindi l'espressione della KCL per il nodo A include i termini $V_A$ e $V_{OUT}$
successivamente, sempre seguendo le istruzioni date si trovano ulteriori
relazioni per i nodi $V_A$ e $V_B$
\begin{equation*}
  sC(V_A - V_B) = \frac{V_B}{2R}
\end{equation*}
assumendo che l'opamp sia ideale e che operi in regime lineare si applica
il principio di corto-circuito virtuale $V_+ = V_-$ da cui
\begin{equation*}
  \begin{gathered}
    \begin{cases}
      \frac{V_{OUT} - V_-}{2R} = \frac{V_-}{\frac65R} \\
      V_B = V_+ = V_-
    \end{cases} \implies
    \frac{V_{OUT} - V_B}{2R} = \frac{V_B}{\frac65R} \\
    V_B = \frac38V_{OUT}
  \end{gathered}
\end{equation*}
ora $V_B$ \'e puramente espressione di $V_{OUT}$ questo permetter\'a
di trovare una relazione analoga per $V_A$
\begin{equation*}
  \begin{gathered}
    \begin{cases}
      sC(V_A - V_B) = \frac{V_B}{2R} \\
      V_B = \frac38V_{OUT}
    \end{cases} \implies
    sC\left(V_A - \frac38V_{OUT}\right) = \frac{\frac38V_{OUT}}{2R} \\
    V_A = \frac38\left(\frac{1 + 2RCs}{2RCs}\right)V_{OUT}
  \end{gathered}
\end{equation*}
ora $V_B$ e $V_A$ sono espressioni unicamente di $V_{OUT}$. Ora si riduce
la prima equazione nodale trovata all'inizio di questa risposta
\begin{equation*}
  V_S = (2 + 2sRC)V_A - sRCV_B -V_{OUT}
\end{equation*}
e si sostituiscono $V_A$ e $V_B$ con le loro espressioni in termini di $V_{OUT}$
\begin{equation*}
  V_S = \left[\frac38(2 + 2sRC)\left(\frac{1 + 2RCs}{2RCs}\right) - 1 - \frac38sRC\right]V_{OUT}
\end{equation*}
l'espressione finale per la funzione di trasferimento complessa
\begin{equation}
  \label{eq:bp-vcvs-tfunc}
  \colorbox{blu}{\boxed{H(s) = \frac{8RCs}{3 + RCs + 3R^2C^2s^2}}}
\end{equation}

\noindent\colorbox{orange}{
  \begin{minipage}{1.0\linewidth}
    \paragraph{Richiesta (c)}
    Dalla funzione di trasferimento dedurre i valori attesi per i parametri del filtro (frequenza
    centrale, guadagno massimo e larghezza banda) e confrontarli con le misure di cui ai
    punti 1a, 1b, 1c.
  \end{minipage}
}
\paragraph{Risposta (c)}
Le parti reali e immaginarie del numeratore e denominatore
\begin{equation*}
  \begin{cases}
    \mathcal{R}_{\omega}^N = 0 \\
    \mathcal{I}_{\omega}^N = 8RC \\
    \mathcal{R}_{\omega}^D = 3 - 3R^2C^2 \\
    \mathcal{I}_{\omega}^D = RC
  \end{cases} \implies
  G(\omega) = \sqrt{\frac{(8RC\omega)^2}{(3 - 3R^2C^2\omega^2)^2 + (RC\omega)^2}}
\end{equation*}
il guadagno \'e ottenuto dall'espressione \ref{eq:da-tfunc-a-gain}
\begin{equation}
  \label{eq:bp-vcvs-gain}
  G(f) = \frac{16\pi RCf}{\sqrt{9 - 68\pi^2R^2C^2f^2 + 144\pi^4R^4C^4f^4}}
\end{equation}
il guadagno massimo pu\'o essere ottenuto semplicemente riducendo l'espressione
del guadagno trovata in precedenza (in funzione della pulsazione $\omega$) nella
seguente forma
\begin{equation*}
  G(\omega) = \frac{8}{\sqrt{\left(\frac{3 - 3R^2C^2\omega^2}{RC\omega}\right)^2 + 1}}
\end{equation*}
il numeratore assume valore minimo per $3 - 3R^2C^2\omega^2 = 0$ e quindi il guadagno \'e
\begin{equation}
  \label{eq:bp-vcvs-gain-max}
  \colorbox{blu}{\boxed{\max_{\omega \in [0, +\infty)} G(\omega) = 8}}
\end{equation}
inoltre sostituendo $\omega \to 2\pi f$ alla condizione di minimizzazione si ricava la
frequenza centrale
\begin{equation}
  \label{eq:bp-vcvs-frequenza-centrale}
  \colorbox{blu}{\boxed{f = \frac{1}{2\pi RC}}}
\end{equation}
Un modo di ottenere la larghezza di banda richiede le frequenze di taglio, si calcolano
vincolando il guadagno quadro alla met\'a del suo guadagno massimo quadro (si usa il guadagno
in funzione della pulsazione per facilitare i calcoli)
\begin{equation*}
  \begin{gathered}
    \frac{8RC\omega}{\sqrt{9 - 17R^2C^2\omega^2 + 9R^4C^4\omega^4}} = \frac{8}{\sqrt2} \\ \implies
    9 - 19R^2C^2\omega^2 + 9R^4C^4\omega^4 = 0
  \end{gathered}
\end{equation*}
un ulteriore sostituzione $R^2C^2\omega^2 \to \xi$ semplifica il problema ulteriormente ad
un polinomio di ordine 2
\begin{equation*}
  \begin{gathered}
    9 - 19\xi + 9\xi^2 = 0 \\
    \xi = \frac{19 \pm \sqrt{37}}{18}
  \end{gathered}
\end{equation*}
sostituendo $\xi \to R^2C^2\omega^2$ e $\omega \to 2\pi f$ si ottiene
\begin{equation*}
  \begin{gathered}
    4\pi^2R^2C^2f^2 = \frac{19 \pm \sqrt{37}}{18} \implies
    f_{L,H} = \frac{1}{2\pi RC}\sqrt{\frac{19 \pm \sqrt{37}}{18}}
  \end{gathered}
\end{equation*}
dunque la larghezza di banda $\Delta f = f_H - f_L$
\begin{equation}
  \label{eq:bp-vcvs-bw}
  \colorbox{blu}{\boxed{\Delta f = \frac{1}{2\pi\sqrt{18}RC}\left[\sqrt{19 + \sqrt{37}} - \sqrt{19 - \sqrt{37}}\right]}}
\end{equation}

\noindent\colorbox{orange}{
  \begin{minipage}{1.0\linewidth}
    \paragraph{Richiesta (d)}
    Dare un’interpretazione qualitativa (e se possibile quantitativa) delle caratteristiche del
    segnale di uscita osservato al punto 1d. Che tipo di segnale appare? A quale ampiezza?
  \end{minipage}
}


\section*{Highpass VCVS Butterworth}
\noindent\colorbox{orange}{
  \begin{minipage}{1.0\linewidth}
    \paragraph{Richiesta (a)}
    Dimostrare, discutendone la risposta ai limiti del dominio delle frequenze reali sulla base
    di semplici argomenti di natura qualitativa, che il circuito proposto svolga effettivamente
    la funzione di filtro passa-alto.
  \end{minipage}
}

\paragraph{Risposta (a)}
\begin{itemize}
\item
  Per $f \to 0$ il condensatore $C_1$ apre il ramo dunque $V_+ = 0$. Se l'opamp
  opera in regime di linearit\'a $V_{OUT} = A_dV_D$
  \begin{equation*}
    V_{OUT} = A_d(V_+ - V_-) = A_d(0 - V_{OUT}) \implies \colorbox{blu}{\boxed{V_{OUT} = \frac{0}{1 + A_d} = 0}}
  \end{equation*}
  nel limite di basse frequenze il circuito ha un andamento compatibile con quello di
  un circuito filtro passa alto.
\item Per $f \to +\infty$ i condensatori $C_1$ e $C_2$ ``corto-circuitano'' i propri rami
  cos\'i che $V_+ = V_s$. Se l'opamp opera in regime di linearit\'a $V_{OUT} = A_dV_D$
  \begin{equation*}
    V_{OUT} = A_d(V_+ - V_-) = A_d(V_s - V_{OUT}) \implies \colorbox{blu}{\boxed{V_{OUT} = \frac{A_d}{1 + A_d}V_s}}
  \end{equation*}
  nel limite di alte frequenze il circuito ha un andamento compatibile con quello di
  un circuito filtro passa alto.
\end{itemize}

\noindent\colorbox{orange}{
  \begin{minipage}{1.0\linewidth}
    \paragraph{Richiesta (b)}
    Nell’ipotesi di idealit\'a dell’operazionale ed assumendo C1 = C2 = C e R1=2, R2=2R,
    derivare analiticamente un’espressione della trasformata di Laplace della funzione di
    trasferimento $A_v(s)=V_{OUT}(s)/V_S(s)$.
  \end{minipage}
}
\paragraph{Risposta (b)}
\begin{itemize}
\item
  Tra $V_-$ e $V_{OUT}$ i rami sono corto-circuiti, di conseguenza $V_- = V_{OUT}$.
\item
  Applico il principio di \textit{cortocircuito virtuale} e di conseguenza $V_+ = V_-$.
\end{itemize}
di conseguenza
\begin{equation*}
  V_B = V_{OUT}
\end{equation*}
le successive equazioni nodali appaiate con il risultato precedente permettono di
ottenere la funzione di trasferimento corretta
\begin{equation*}
  \begin{gathered}
    sC(V_A - V_B) = \frac{V_{OUT}}{2R} \iff V_A = \left(\frac{1 + 2sRC}{2sRC}\right)V_{OUT} \\
    sC(V_S - V_A) = sC(V_A - V_B) + \frac{V_A - V_{OUT}}{R} \\
    sRC(V_S - V_A) = sRC(V_A - V_B) + V_A - V_{OUT} \\
    sRCV_S = (1 + 2sRC)V_A - (1 + sRC)V_{OUT} \\
    sRCV_S = \left[(1 + 2sRC)\left(\frac{1 + 2sRC}{2sRC}\right) - (1 + sRC)\right]V_{OUT} \\
    2s^2R^2C^2V_S = \left[(1 + 2sRC)(1 + 2sRC) - 2sRC(1 + sRC)\right]V_{OUT} \\
  \end{gathered}
\end{equation*}
da cui
\begin{equation}
  \label{eq:hp-vcvs-butterworth-tfunc}
  \colorbox{blu}{\boxed{H(s) = \frac{2s^2R^2C^2}{1 + 2RCs + 2R^2C^2s^2}}}
\end{equation}

\noindent\colorbox{orange}{
  \begin{minipage}{1.0\linewidth}
    \paragraph{Richiesta (c)}
    Descrivere le propriet\'a analitiche di $A_v(s)$ (ovvero calcolarne zeri e poli) e da queste
    dedurre la stabilit\'a del filtro.
  \end{minipage}
}
\paragraph{Risposta (c)}
Per studiare gli zero della funzione di trasferimento complessa si studiano
le radici del numeratore della forma canonica.
\begin{equation*}
  2s^2R^2C^2 = 0
\end{equation*}
In questo caso \'e presente uno zero del secondo ordine
\begin{equation}
  \label{eq:hp-vcvs-butterworth-zeros}
  \colorbox{blu}{\boxed{z^2 = 0}}
\end{equation}
Per studiare i poli della funzione di trasferimento complessa si studiano
le radici del denominatore della forma canonica.
\begin{equation*}
  \begin{cases}
    \xi \equiv 2RCs \\
    1 + 2\xi + 2\xi^2 = 0
  \end{cases} \implies
  \xi_{1,2} = \frac{-2 \pm \sqrt{4-4\cdot 1\cdot 2}}{2\cdot 2} = -\frac12 \pm j\frac12
\end{equation*}
In questo sono presenti due poli del secondo ordine
\begin{equation}
  \label{eq:hp-vcvs-butterworth-poles}
  \colorbox{blu}{\boxed{p_{1,2} = -\frac{1}{4RC} \pm j\frac{1}{4RC}}}
\end{equation}
il filtro \'e \textbf{stabile} perch\'e entrambi i suoi poli hanno parte reale negativa. \\

\noindent\colorbox{orange}{
  \begin{minipage}{1.0\linewidth}
    \paragraph{Richiesta (d)}
    Dalla funzione di trasferimento determinare i valori attesi per le caratteristiche del filtro
    (guadagno massimo, frequenza di taglio, andamento di modulo/fase in funzione della
    frequenza) e confrontarli con quanto misurato al punto 1c.
  \end{minipage}
}
\paragraph{Risposta (d)}
Il guadagno si ottiene mediante \ref{eq:da-tfunc-a-gain}, calcolando le parti reali ed immaginarie del
numeratore e denominatore della forma canonica della funzione di trasferimento $H(s)$ si ottiene
\begin{equation*}
  \begin{gathered}
    \begin{cases}
      \mathcal{R}_{\omega}^N = -2R^2C^2\omega^2 \\
      \mathcal{I}_{\omega}^N = 0 \\
      \mathcal{R}_{\omega}^D = 1-2R^2C^2\omega^2 \\
      \mathcal{I}_{\omega}^D = 2RC\omega
    \end{cases} \implies
    G(\omega) = \sqrt{\frac{(-2R^2C^2\omega^2)^2}{(1-2R^2C^2\omega^2)^2 + (2RC\omega)^2}}
  \end{gathered}
\end{equation*}
semplificando l'espressione e sostituendo $\omega \to 2\pi f$ si ottiene il guadagno in funzione
della frequenza
\begin{equation}
  \label{eq:hp-vcvs-butterworth-gain}
  \colorbox{blu}{\boxed{G(f) = \frac{8\pi^2R^2C^2f^2}{\sqrt{1+64\pi^4R^4C^4f^4}}}}
\end{equation}
Per calcolare l'espressione della fase bisogna prima identificare il semi-piano
di appartenenza, il segno della parte reale \'e determinata da
\begin{equation*}
  (-2R^2C^2\omega^2)(1-2R^2C^2\omega^2)
\end{equation*}
il segno non \'e ben determinato per ogni $\omega \in [0, +\infty)$. Invece, il segno
della parte immaginaria \'e determinata da
\begin{equation*}
  -(-2R^2C^2\omega^2)(2RC\omega) = 4R^3C^3\omega^3 > 0
\end{equation*}
di conseguenza il semipiano \'e quello \textit{immaginario positivo} da cui
la fase $\phi$ in funzione della frequenza viene calcolata attraverso la \ref{eq:da-tfunc-a-fase-immaginario-positivo}
e la sostituzione $\omega \to 2\pi f$
\begin{equation}
  \label{eq:hp-vcvs-butterworth-fase}
  \colorbox{blu}{\boxed{\phi(f) = \frac{\pi}{2} + \arctan\left(\frac{1-8\pi^2R^2C^2f^2}{4\pi RCf}\right)}}
\end{equation}
Il numeratore e denominatore del guadagno sono di segno concorde (naturalmente), il numeratore
\'e monotono cresecente e il denominatore monotono decrescente, di conseguenza il guadagno \'e
monotono crescente di conseguenza il massimo si trova nel limite superiore del dominio
\begin{equation}
  \label{eq:hp-vcvs-butterworth-gain-max}
  \colorbox{blu}{\boxed{\max_{f \in [0, +\infty)}G(f) = \lim_{f\to +\infty}G(f) = 1}}
\end{equation}
la frequenza di taglio \'e per definizione la frequenza per cui il guadagno quadro \'e
met\'a del massimo (appena trovato)
\begin{equation*}
  \frac{8\pi^2R^2C^2f^2}{\sqrt{1+64\pi^4R^4C^4f^4}} = \frac{1}{\sqrt2} \implies
  2(8\pi^2R^2C^2f^2)^2 = 1+64\pi^4R^4C^4f^4
\end{equation*}
\begin{equation}
  \label{eq:hp-vcvs-butterworth-fcut}
  \colorbox{blu}{\boxed{f = \frac{1}{2\sqrt{2}\pi RC}}}
\end{equation}

\paragraph{Richiesta (e)}
Dare un’interpretazione qualitativa (e se possibile quantitativa) delle caratteristiche del
segnale di uscita osservato ai punti 1d, 1e.


\section*{Lowpass VCVS Butterworth}

\noindent\colorbox{orange}{
  \begin{minipage}{1.0\linewidth}
    \paragraph{Richiesta (a)}
    Dimostrare, sulla base di semplici argomenti di natura qualitativa, che il circuito proposto
    svolga effettivamente la funzione di filtro passa-basso, discutendone la risposta ai limiti
    del dominio delle frequenze reali.
  \end{minipage}
}

\paragraph{Risposta (a)}
\begin{itemize}
\item Per $f \to 0$
  i condensatori agiscono da rami aperti, le resistenze hanno impedenze infinite che vengono
  assorbite dall'impedenze ideali all'ingresso dell'opamp. Questa \'e la configurazione di
  un voltage follower. La tensione d'uscita \'e la tensione d'ingresso, il guadagno \'e unitario.
  
\item Per $f \to \infty$
  i condensatori agiscono da cortocircuiti, l'ingresso non-invertente viene messo a terra,
  l'ingresso invertente nell'ipotesi di corto-circuito virtuale deve eguagliare la tensione,
  la tensione di uscita deve essere nulla.
\end{itemize}


\noindent\colorbox{orange}{
  \begin{minipage}{1.0\linewidth}
    \paragraph{Richiesta (b)}
    Derivare analiticamente la funzione di trasferimento complessa $A_v(s)=V_{OUT}(s)/V_S(s)$
    nell’ipotesi che l’operazionale sia ideale (suggerimento: si determini la tensione nel nodo
    A del circuito in termini di $V_S$ e $V_{OUT}$; quindi si esprima $V_A$ in termini di $V_{OUT}$ assumendo
    che l'opAmp funzioni in regime lineare ed utilizzando il principio di cortocircuito
    virtuale; si eguaglino infine le espressioni ottenute per ottenere la funzione di
    trasferimento). Per semplicit\'a di calcolo, si assuma inoltre (ipotesi compatibile con i
    valori e le incertezze nominali dei componenti) che R1=R2=R, C1=C, C2=2C.
  \end{minipage}
}

\paragraph{Risposta (b)}
Ipotesi: l'operazionale \'e ideale ed opera in regime lineare.
\begin{itemize}
\item
  Tra $V_-$ e $V_{OUT}$ i rami sono corto-circuiti, di conseguenza $V_- = V_{OUT}$.
\item
  Applico il principio di \textit{cortocircuito virtuale} e di conseguenza $V_+ = V_-$.
\end{itemize}
di conseguenza
\begin{equation*}
  V_B = V_{OUT}
\end{equation*}
le successive equazioni nodali appaiate con il risultato precedente permettono di
ottenere la funzione di trasferimento corretta
\begin{equation*}
  \begin{gathered}
    \frac{V_A - V_B}{R} = sCV_{OUT} \iff V_A = (1 + sRC)V_{OUT} \\
    \frac{V_S - V_A}{R} = \frac{V_A - V_B}{R} + 2sC(V_A - V_{OUT}) \\
    V_S - V_A = V_A - V_B + 2sRC(V_A - V_{OUT}) \\
    V_S = (2 + 2sRC)V_A - (1 + 2sRC)V_{OUT} \\
    V_S = \left[(2 + 2sRC)(1 + sRC) - (1 + 2sRC)\right]V_{OUT} \\
  \end{gathered}
\end{equation*}
da cui
\begin{equation}
  \label{eq:lp-vcvs-butterworth-tfunc}
  \colorbox{blu}{\boxed{H(s) = \frac{1}{1 + 2RCs + 2R^2C^2s^2}}}
\end{equation}

\noindent\colorbox{orange}{
  \begin{minipage}{1.0\linewidth}
    \paragraph{Richiesta (c)}
    Dalla funzione di trasferimento dedurre le caratteristiche del filtro osservate al punto 1)
    (guadagno a centro-banda, frequenza di taglio, andamento di modulo/fase in funzione
    della frequenza) e confrontarle con quanto misurato ai punti 1c, 1d, 1e.
  \end{minipage}
}

\paragraph{Risposta (c)}
il guadagno \'e ottenuto per mezzo dell'esspressione \ref{eq:da-tfunc-a-gain}
\begin{equation*}
  \begin{gathered}
    \begin{cases}
      \mathcal{R}_{\omega}^N = 1 & \mathcal{I}_{\omega}^N = 0 \\
      \mathcal{R}_{\omega}^D = 1 - 2R^2C^2\omega^2 & \mathcal{I}_{\omega}^D = -2RC\omega
    \end{cases} \\
    G(\omega) = \frac{1}{\sqrt{(1 - 2R^2C^2\omega^2)^2 + (-2RC\omega)^2}} = \frac{1}{\sqrt{1 + (4R^4C^4)\omega^4}}
  \end{gathered}
\end{equation*}
\begin{equation}
  \label{eq:lp-vcvs-butterworth-gain}
  \colorbox{blu}{\boxed{G(f) = \frac{1}{\sqrt{1 + (64\pi^4R^4C^4)f^4}}}}
\end{equation}
per determinare la fase si usa \ref{eq:da-tfunc-a-fase-immaginario-negativo}
\begin{equation*}
  \phi(\omega) = \arctan\left(\frac{1 - 2R^2C^2\omega^2}{2RC\omega}\right)
\end{equation*}
\begin{equation}
  \label{eq:lp-vcvs-butterworth-phase}
  \colorbox{blu}{\boxed{\phi(f) = -\frac{\pi}{2} + \arctan\left(\frac{1 - 8\pi^2R^2C^2f^2}{4\pi RCf}\right)}}
\end{equation}


Frequenza di taglio.
\begin{equation*}
  \begin{gathered}
    G(f) = \frac{\max\{G(\omega)\}}{\sqrt2} = \frac{1}{\sqrt2} \implies \frac{1}{\sqrt{1 + (64\pi^4R^4C^4)f^4}} = \frac{1}{\sqrt2} \\
  \end{gathered}
\end{equation*}
\begin{equation}
  \label{eq:lp-vcvs-butterworth-fcut}
  \colorbox{blu}{\boxed{f = \frac{1}{2\pi\sqrt2\cdot RC}}}
\end{equation}

\paragraph{Richiesta (d)}
Dare un’interpretazione qualitativa (e se possibile quantitativa) delle caratteristiche del
segnale di uscita osservato al punto 1f. Che tipo di segnale appare? A quale ampiezza?


\section*{Lowpass VCVS}
\noindent\colorbox{orange}{
  \begin{minipage}{1.0\linewidth}
    \paragraph{Richiesta (a)}
    Dimostrare, sulla base di semplici argomenti di natura qualitativa (ovvero prescindendo
    dall’espressione analitica della funzione di trasferimento) che il circuito proposto svolga
    effettivamente la funzione di filtro passa-basso, discutendone la risposta ai limiti del dominio
    delle frequenze reali e verificare il limite del guadagno a bassa frequenza ottenuto al punto
    1c.
  \end{minipage}
}

\paragraph{Risposta (a)}
\begin{itemize}
\item
  Per $f \to 0$ i condensatori $C_1$ e $C_2$ aprono i loro rami. Il circuito opera
  come un amplificatore \underline{non}-invertente. La tensione al nodo A \'e la
  tensione sorgente \textit{partita} dal partitore di tensione $(R_1, R_2)$
  \begin{equation*}
    V_A = \frac12 V_s
  \end{equation*}
  $R_3$ \'e una resistenza finita trascurabile rispetto all'impendenza ``grande'' dell'opamp,
  di conseguenza la tensione dell'ingresso invertente
  \begin{equation*}
    V_+ = \frac12 V_s
  \end{equation*}
  In maniera del tutto analoga a quanto detto in precedenza, la tensione dell'ingresso
  invertente
  \begin{equation*}
    V_- = \frac12 V_{OUT}
  \end{equation*}
  Se l'opamp opera in regime di linearit\'a $V_{OUT} = A_dV_d$
  \begin{equation*}
    V_{OUT} = \frac12 A_d(V_s - V_{OUT}) \implies \colorbox{blu}{\boxed{V_{OUT} = \frac{A_d}{2 + A_d}V_s}}
  \end{equation*}
  nel limite di basse frequenze il circuito ha un andamento compatibile con quello di
  un circuito filtro passa basso.
\item Per $f \to +\infty$ il condensatore $C_1$ ``corto-circuita'' il proprio ramo
  cos\'i che $V_+ = 0$. Se l'opamp opera in regime di linearit\'a $V_{OUT} = A_dV_d$
  \begin{equation*}
    V_{OUT} = A_d(0 - V_-) = -A_dV_{OUT} \implies \colorbox{blu}{\boxed{V_{OUT} = \frac{0}{1 + A_d} = 0}}
  \end{equation*}
  nel limite di alte frequenze il circuito ha un andamento compatibile con quello di
  un circuito filtro passa basso.
\end{itemize}

\noindent\colorbox{orange}{
  \begin{minipage}{1.0\linewidth}
    \paragraph{Richiesta (b)}
    Derivare analiticamente la trasformata di Laplace della funzione di trasferimento $A_v(s) =
    V_{OUT}(s)/V_S(s)$ nell’ipotesi che l’operazionale sia ideale ed i valori di resistenze e capacit\'a
    siano pari a quelli nominali, ovvero valgano le relazioni C1=C2=C, R1=R2=R, R3=2R,
    R4=R5 (\textit{suggerimento: si determini la tensione nel punto A del circuito in termini di VS e
      $V_{OUT}$; si determini quindi la tensione all’ingresso non invertente dell’operazionale e
      l’ulteriore relazione tra $V_A$ e $V_{OUT}$ assumendo che l'operazionale funzioni in regime lineare}).
  \end{minipage}
}

\paragraph{Risposta (b)}
La richiesta porta con se un suggerimento, seguendolo per il nodo A vale
\begin{equation*}
  \frac{V_S - V_A}{R} = \frac{V_A}{R} + \frac{V_A - V_B}{2R} + sC(V_A - V_{OUT})
\end{equation*}
quindi l'espressione della KCL per il nodo A include i termini $V_A$ e $V_{OUT}$
successivamente, sempre seguendo le istruzioni date si trovano ulteriori
relazioni per i nodi $V_A$ e $V_B$
\begin{equation*}
  \frac{V_A - V_B}{2R} = sCV_B
\end{equation*}
assumendo che l'opamp sia ideale e che operi in regime lineare si applica
il principio di corto-circuito virtuale $V_+ = V_-$ da cui
\begin{equation*}
  \begin{gathered}
    \begin{cases}
      \frac{V_{OUT} - V_-}{R} = \frac{V_-}{R} \\
      V_B = V_+ = V_-
    \end{cases} \implies
    \frac{V_{OUT} - V_B}{R} = \frac{V_B}{R} \\
    V_B = \frac12V_{OUT}
  \end{gathered}
\end{equation*}
ora $V_B$ \'e puramente espressione di $V_{OUT}$ questo permetter\'a
di trovare una relazione analoga per $V_A$
\begin{equation*}
  \begin{gathered}
    \begin{cases}
      \frac{V_A - V_B}{2R} = sCV_B \\
      V_B = \frac12V_{OUT}
    \end{cases} \implies
    \frac{V_A - \frac12V_{OUT}}{2R} = sC\frac12V_{OUT} \\
    V_A = \frac12(1 + 2sRC)V_{OUT}
  \end{gathered}
\end{equation*}
ora $V_B$ e $V_A$ sono espressioni unicamente di $V_{OUT}$. Ora si riduce
la prima equazione nodale trovata all'inizio di questa risposta
\begin{equation*}
  V_S = \frac12(5 + 2sRC)V_A - \frac12V_B - sRCV_{OUT}
\end{equation*}
e si sostituiscono $V_A$ e $V_B$ con le loro espressioni in termini di $V_{OUT}$
\begin{equation*}
  V_S = \frac14\left[(5 + 2sRC)(1 + 2sRC) - 1 - 4sRC\right]V_{OUT}
\end{equation*}
l'espressione finale per la funzione di trasferimento complessa
\begin{equation}
  \label{eq:lp-vcvs-tfunc}
  \colorbox{blu}{\boxed{H(s) = \frac{1}{(1 + RCs)^2}}}
\end{equation}

\noindent\colorbox{orange}{
  \begin{minipage}{1.0\linewidth}
    \paragraph{Richiesta (c)}
    Dalla funzione di trasferimento dedurre le caratteristiche del filtro osservate al punto 1)
    (guadagno a centro-banda, frequenza naturale, andamento di modulo/fase in funzione della
    frequenza) e confrontarli con quanto misurato ai punti 1c, 1d, 1e.
  \end{minipage}
}
\paragraph{Risposta (c)}
Le parti reali e immaginarie del numeratore e denominatore
\begin{equation*}
  \begin{cases}
    \mathcal{R}_{\omega}^N = 1 \\
    \mathcal{I}_{\omega}^N = 0 \\
    \mathcal{R}_{\omega}^D = 1 - R^2C^2\omega^2 \\
    \mathcal{I}_{\omega}^D = 2RC\omega
  \end{cases} \implies
  G(\omega) = \sqrt{\frac{1}{(1 - R^2C^2\omega^2)^2 + (2RC\omega)^2}}
\end{equation*}
il guadagno \'e ottenuto dall'espressione \ref{eq:da-tfunc-a-gain}
\begin{equation}
  \label{eq:lp-vcvs-gain}
  G(f) = \frac{1}{1 + 4\pi^2R^2C^2f^2}
\end{equation}
il denominatore \'e positivo monotono crescente, limitato inferiormente a $G(f) = 1$,
di conseguenza il guadagno massimo, essando il reciproco del denominatore \'e proprio
quell'estremo inferiore
\begin{equation}
  \label{eq:lp-vcvs-gain-max}
  \colorbox{blu}{\boxed{\max_{f \in [0, +\infty)} G(f) = G(0) = 1}}
\end{equation}
la frequenza centrale \'e la frequenza per cui il guadagno \'e massimo
\begin{equation}
  \label{eq:lp-vcvs-frequenza-centrale}
  \colorbox{blu}{\boxed{f = 0}}
\end{equation}
ad ogni polo della funzione di trasferimento complessa $H(s)$ corrisponde una
frequenza naturale\footnotemark

\footnotetext{
  Il viceversa non \'e vero in generale, alcune frequenze naturali
  corrispondono a particolari condizioni iniziali
  ``Desoer, Charles (1969). Basic circuit theory. McGraw-Hill pag. 643''
}

\begin{equation}
  \label{eq:lp-vcvs-frequenza-naturale}
  \colorbox{blu}{\boxed{f = \frac{1}{2\pi RC}}}
\end{equation}

Per calcolare l'espressione della fase bisogna prima identificare il semi-piano di appartenenza,
il sengo della parte reale \'e determinata da
\begin{equation*}
  1-R^2C^2\omega^2
\end{equation*}
il segno non \'e be determinato per ogni $\omega \in [0, +\infty)$. Invece, il segno della parte
immaginaria \'e determinata da
\begin{equation*}
  2RC\omega > 0
\end{equation*}
di conseguenza il semipiano \'e quello \textit{immaginario positivo} da cui la fase $\phi$ in
funzione della frequenza viene calcolata attraverso la \ref{eq:da-tfunc-a-fase-immaginario-positivo}
e la sostituzione $\omega \to 2\pi f$
\begin{equation}
  \label{eq:lp-vcvs-fase}
  \colorbox{blu}{\boxed{\phi(f) = -\frac{\pi}{2} + \arctan\left(\frac{1 - 4\pi^2R^2C^2f^2}{4\pi RCf}\right)}}
\end{equation}

\noindent\colorbox{orange}{
  \begin{minipage}{1.0\linewidth}
    \paragraph{Richiesta (d)}
    Dare un’interpretazione qualitativa e quantitativa delle caratteristiche del segnale di uscita
    osservato al punto 1f. Che tipo di segnale appare? A quale ampiezza?
  \end{minipage}
}


\section*{Notch 2}
\noindent\colorbox{orange}{
  \begin{minipage}{1.0\linewidth}
    \paragraph{Richiesta (a)}
    Utilizzando semplici argomenti di natura qualitativa (ovvero senza effettuare il calcolo
    della funzione di trasferimento, come richiesto al punto successivo) mostrare che il
    circuito svolga effettivamente la funzione di notch, discutendone la risposta ai limiti del
    dominio delle frequenze reali e verificando che i limiti asintotici del guadagno (sia per
    $f \to 0$ che per $f \to \infty$) coincidano con quanto osservato al punto 1b.
  \end{minipage}
}
\paragraph{Risposta (a)}
Qualitativamente il circuito si tratta di un amplificatore pilotato
dall'ingresso non invertente che viene preceduto da un nodo che
ridirige il flusso di corrente ad un blocco secondario da studiare
qualitativamente. Nei limiti di frequenze il sotto-circuito agisce
come se non esistesse. Il ramo in retroazione dell'ingresso invertente
agisce come un \textit{dimezzatore di tensione}, dunque se l'impedenza
effettiva del sotto-circuito equivale alla configurazione all'ingresso
invertente il guadagno differenziale sar\'a nullo e il filtro agira
come \textit{notch-filter}. Con uno studio quantitativo si potra calcolare
quanti notch sono presenti e dove sono disposti in funzione della frequenza.

\noindent\colorbox{orange}{
  \begin{minipage}{1.0\linewidth}
    \paragraph{Richiesta (b)}
    Derivare analiticamente la funzione di trasferimento (in trasformata di Laplace) A(s)
    nell'ipotesi di linearità (da giustificare a posteriori, vedi punto d) ed assumendo per
    semplicità che R1=R2=R', R3=R=20kW, R4=R5=R/2, C1=C=1nF e C2=100 C.
    Si suggerisce di procedere:
    \begin{itemize}
    \item[i]
      ottenendo un'espressione dell'impedenza di ingresso ZIN della rete inserita nel
      riquadro tratteggiato nello schema circuitale, ovvero quella vista dall'ingresso non
      invertente dell'operazionale U1 (\textit{suggerimento: si determini un'espressione della
        tensione all'ingresso non invertente dell'operazionale U1 in termini della corrente I1
        indicata in figura e considerando la relazione tra quella corrente e I2});
    \item[ii]
      ricavando una relazione tra VOUT, VS, R e ZIN sulla base dello schema equivalente
      riportato in calce ed infine calcolando il rapporto VOUT/VS in funzione di s (o, più
      convenientemente, della variabile complessa adimensionale $\xi=sRC$).
    \end{itemize}
  \end{minipage}
}
\paragraph{Risposta (b.i)}
Seguendo il suggerimento
\begin{equation*}
  I_1 \equiv sC(V_{+U1} - V_{+U2})
\end{equation*}
nell'ipotesi di op-amp ideale, +U2 non assorbe corrente, di conseguenza
\begin{equation*}
  I_1 = \frac{V_{+U2} - V_3}{R/2}
\end{equation*}
inoltre lo schema circuitale dell'op-amp U2 \'e quello di un inseguitore
di tensione, di conseguenza
\begin{equation*}
  V_{+U2} = V_{U2} \implies I_1 = \frac{V_{U2} - V_3}{R/2}
\end{equation*}
la corrente all'uscita dell'op-amp U2
\begin{equation*}
  I_2 \equiv 100sC(V_{U2} - V_3) \implies I_2 = 100sC\left(\frac{R}{2}I_1\right)
\end{equation*}
la KCL al nodo 3
\begin{equation*}
  \begin{gathered}
    I_1 + I_2 = I_3 \equiv \frac{V_3}{R/2} \\
    \implies V_3 = \frac{R}{2}(1 + 50sRC)I_1
  \end{gathered}
\end{equation*}
dalla seconda espressione di $I_1$
\begin{equation*}
  \begin{gathered}
    I_1 = \frac{V_{+U2} - \frac{R}{2}(1 + 50sRC)I_1}{R/2} \\
    \implies V_{+U2} = \frac{R}{2}(2 + 50sRC)I_1
  \end{gathered}
\end{equation*}
dalla prima espressione/definizione di $I_1$
\begin{equation*}
  \begin{gathered}
    I_1 = sC(V_{+U1} - R(1 + 25sRC)I_1) \\ \implies
    (1 + sRC + 25sR^2C^2)I_1= sCV_{+U1}
  \end{gathered}
\end{equation*}
l'impedenza d'ingresso $Z_{in}$ si ottiene dal rapporto fra la tensione
d'ingresso $V_{+U1}$ e la corrente d'ingresso $I_1$
\begin{equation}
  \label{eq:notch2-zin}
  \colorbox{blu}{\boxed{Z_{in} \equiv \frac{V_{+U1}}{I_1} = \frac{1}{sC}(1 + sRC + 25sR^2C^2)}}
\end{equation}

\paragraph{Risposta (b.ii)}
La tensione all'ingresso invertente $V_-$ \'e met\'a della tensione sorgente $V_s$
(partitore di tensione con resistenze uguali dimezzano la tensione)
\begin{equation*}
  V_+ = \frac12V_s + \frac12V_{OUT}
\end{equation*}
la tensione all'ingresso \underline{non}-invertente risulta sempre come risultato
di un partiore composto da $R$ e $Z_{in}$
\begin{equation*}
  V_- = \frac{Z_{in}}{R + Z_{in}}V_s
\end{equation*}
nell'ipotesi di op-amp ideale che opera in regime di linearit\'a si applica il principio
di corto-circuito virtuale
\begin{equation*}
  V_+ = V_- \implies \frac12V_s + \frac12V_{OUT} = \frac{Z_{in}}{R + Z_{in}}V_s
\end{equation*}
la forma generale della funzione di trasferimento parametrizzata da $Z_{in}$
\begin{equation*}
  H(s) = \frac{Z_{in} - R}{Z_{in} + R}
\end{equation*}
espandendo l'espressione di $Z_{in}$ si ottiene
\begin{equation}
  \label{eq:notch2-tfunc}
  \colorbox{blu}{\boxed{H(s) = \frac{1 + 25\xi^2}{1 + 2\xi + 25\xi^2} \qquad \xi \equiv sRC}}
\end{equation}

\noindent\colorbox{orange}{
  \begin{minipage}{1.0\linewidth}
    \paragraph{Richiesta (c)}
    Verificare che l'espressione ottenuta per la funzione di trasferimento soddisfi la
    condizione necessaria a garantire la stabilità del filtro in zona lineare.
  \end{minipage}
}
\paragraph{Risposta (c)}
La stabilit\'a del filtro si verifica studiando i poli della funzione di
trasferimento complessa $H(s)$ (ovvero gli zeri del denominatore)
\begin{equation*}
  \begin{gathered}
    1 + 2\xi + 25\xi^2 = 0 \implies \xi = \frac{-2 \pm i\sqrt{96}}{50} \\ \implies
    \colorbox{blu}{\boxed{\mathcal{R}\{\xi\} = -\frac{1}{25} < 0}}
  \end{gathered}
\end{equation*}
la parte reale dei poli \'e negativa, per il terzo criterio di stabilit\'a
questo implica che il filtro \'e stabile e lineare. \\

\noindent\colorbox{orange}{
  \begin{minipage}{1.0\linewidth}
    \paragraph{Richiesta (d)}
    Dedurre i valori attesi (senza incertezze) per i parametri del filtro (frequenza centrale,
    larghezza di banda e fattore di merito) e confrontarli con le misure di cui al punto 1b (si
    consiglia di esprimere la funzione di trasferimento sul dominio delle frequenze reali in
    termini della grandezza adimensionale $u = \omega RC$).
  \end{minipage}
}
\paragraph{Risposta (d)}
Seguendo il consiglio della richiesta per la scelta di parametri, le parti reali
e immaginarie del numeratore e denominatore della funzione di trasferimento
complessa in forma canonica
\begin{equation*}
  \begin{cases}
    \mathcal{R}_{\omega}^N = 1 - 25u^2 \\
    \mathcal{I}_{\omega}^N = 0 \\
    \mathcal{R}_{\omega}^D = 1 - 25u^2 \\
    \mathcal{I}_{\omega}^D = 2u
  \end{cases}
\end{equation*}
il guadagno \'e ottenuto dall'espressione \ref{eq:da-tfunc-a-gain}
\begin{equation*}
  G(u) = \frac{1 - 25u^2}{\sqrt{(1 - 25u^2)^2 + (2u)^2}}
\end{equation*}
nel contesto attuale, ovvero nello studio di un filtro notch, la frequenza
centrale \'e quella frequenza per cui il guadagno si annulla (il \textit{notch})
\begin{equation*}
  1 - 25u^2 = 0
\end{equation*}
sostituendo $u \to \omega RC$ e $\omega \to 2\pi f$ si ottiene la frequenza centrale
\begin{equation}
  \label{eq:notch2-frequenza-centrale}
  \colorbox{blu}{\boxed{f = \frac{1}{10\pi RC}}}
\end{equation}
per studiare la larghezza di banda servono le frequenze di taglio e di conseguenza
anche il guadagno massimo, ri-arrangiando l'espressione del guadagno in termini di $u$
(nell'ipotesi che il numeratore non sia nullo quindi diverso dalla frequenza centrale)
si ottiene
\begin{equation*}
  G(u) = \frac{1}{\sqrt{1 + \left(\frac{2u}{1 - 25u^2}\right)^2}}
\end{equation*}
la funzione di trasferimento \'e il reciproco della radice di un espressione
quadratica con minimo uguale ad $1$ per valori reali, dunque il guadagno massimo
\begin{equation*}
  \max_{u \in [0, +\infty)} G(u) = 1
\end{equation*}
ora le frequenze di taglio si determinano come quelle frequenze per cui il guadagno
quadro \'e uguale alla met\'a del suo massimo quadro
\begin{equation*}
  \begin{gathered}
    \frac{1 - 25u^2}{\sqrt{(1 - 25u^2)^2 + (2u)^2}} = \frac{1}{\sqrt2} \\ \implies
    (1 - 25u^2)^2 = (2u)^2 \implies
    \begin{cases}
      1 - 25u_+^2 = +2u_+ \\
      1 - 25u_-^2 = -2u_-
    \end{cases} \\ \implies
    u_+ = \frac{-1 \pm \sqrt{26}}{25} \quad \wedge \quad u_- = \frac{1 \pm \sqrt{26}}{25}
  \end{gathered}
\end{equation*}
una frequenza ammissibile deve essere positiva, di conseguenza
\begin{equation*}
  u_{L,H} = \frac{\sqrt{26} \pm 1}{25}
\end{equation*}
attraverso la sostituzione $u \to \omega RC$ e $\omega \to 2\pi f$ si ottengono
le frequenze di taglio
\begin{equation*}
  f_{L,H} = \frac{\sqrt{26} \pm 1}{50\pi RC}
\end{equation*}
dunque la larghezza di banda
\begin{equation}
  \label{eq:notch2-bw}
  \colorbox{blu}{\boxed{\Delta f = \frac{1}{25\pi RC}}}
\end{equation}
infine il fattore di merito (Q-factor) si ottiene dal rapporto fra \ref{eq:notch2-frequenza-centrale}
e \ref{eq:notch2-bw}
\begin{equation}
  \label{eq:notch2-q-factor}
  \colorbox{blu}{\boxed{Q = \frac{25\pi RC}{10\pi RC} = \frac52}}
\end{equation}
questo valore \'e indipendente dai componenti del circuito. \\

\noindent\colorbox{orange}{
  \begin{minipage}{1.0\linewidth}
    \paragraph{Richiesta (e)}
    Discutere gli andamenti osservati nei plot di Bode di cui al punto 1a e
    giustificarli alla luce delle proprietà di simmetria della funzione di trasferimento.
  \end{minipage}
} \smallskip \\
\paragraph{Risposta (e)}

\noindent\colorbox{orange}{
  \begin{minipage}{1.0\linewidth}
    \paragraph{Richiesta (f)}
    Inviata in VS un'onda quadra di frequenza pari ad f0 ed ampiezza 500mV,
    osservare all'oscilloscopio ingresso (C1), uscita (C2) e la loro differenza (a tale scopo
    occorre aggiungere un ulteriore canale mediante il menù a tendina ``Add Channel →
    Math → Simple (software)'' associando la differenza dei due canali) e discutere
    qualitativamente e quantitativamente le caratteristiche di quest'ultima.
  \end{minipage}
} \smallskip \\
\paragraph{Risposta (f)}


\section*{Notch Bandpass Sottrattore}

\noindent\colorbox{orange}{
    \begin{minipage}{1.0\linewidth}
      \paragraph{Richiesta (a)}
      Sulla base di semplici argomenti di natura qualitativa (ovvero senza effettuare il calcolo
      della funzione di trasferimento, come richiesto al punto successivo), mostrare che il
      blocco 1 del circuito svolga effettivamente la funzione di filtro passa-banda,
      discutendone la risposta ai limiti del dominio delle frequenze reali.
    \end{minipage}
}
\paragraph{Risposta (a)}
\begin{itemize}
\item
  Per $f \to 0$ i rami dei condensatori $C_1$ e $C_2$ si aprono, per impedenza d'ingresso
  ``elevata'' la resistenza $R_3$ \'e trascurabile per cui
  \begin{equation*}
    V_- = V_{BP}
  \end{equation*}
  e nella regione di operazione lineare dell'op-amp
  \begin{equation*}
    V_{BP} = A_d(0 - V_{BP}) \implies \colorbox{blu}{\boxed{V_{BP} = \frac{0}{1 + A_d} = 0}}
  \end{equation*}
\item
  Per $f \to \infty$ il ramo del condensatore $C_2$ si ``corto-circuita'', la tensione
  al nodo A, di conseguenza la tensione nell'ingresso invertente dell'op-amp
  \begin{equation*}
    V_- = V_{BP}
  \end{equation*}
  nella regione di operazione lineare dell'op-amp
  \begin{equation*}
    V_{BP} = A_d(0 - V_{BP}) \implies \colorbox{blu}{\boxed{V_{BP} = \frac{0}{1 + A_d} = 0}}
  \end{equation*}

\end{itemize}

\noindent\colorbox{orange}{
  \begin{minipage}{1.0\linewidth}
    \paragraph{Richiesta (b)}
    Derivarne analiticamente la funzione di trasferimento $A_{BP}(s)=V_{BP}(s)/V_S(s)$ nell'ipotesi
    che l'operazionale sia ideale ed assumendo per semplicit\'a che $R_1=R_2=R=10 kW$, $R3=2R$ e
    $C_1=C_2=C=10 nF$ \textit{
      (suggerimento: si determini la tensione nel punto A del circuito in termini
      di $V_S$ e $V_{BP}$, assumendo che l'operazionale operi in regime lineare come risulta dalle
      precedenti misure; indipendentemente si esprima $V_A$ in termini di $V_{BP}$; si eguaglino
      infine le espressioni ottenute per ottenere la funzione di trasferimento).
    }
  \end{minipage}
}
\paragraph{Risposta (b)}
\begin{equation*}
  \frac{V_S - V_A}{R} = \frac{V_A}{R} + sC(V_A - V_-) + sC(V_A - V_{BP})
\end{equation*}

\begin{equation*}
  sC(V_A - V_-) = \frac{V_- - V_{BP}}{2R}
\end{equation*}
per il principio di corto-circuito virtuale, siccome l'ingresso \underline{non}-invertente viene messo
a terra $V_+ = 0$
\begin{equation*}
  \begin{gathered}
    V_+ = V_- \implies sCV_A = -\frac{V_{BP}}{2R} \\ \implies
    V_A = -\frac{V_{BP}}{2sRC}
  \end{gathered}
\end{equation*}
ricordandoci delle conseguenze del principio di corto-circuito virtuale, l'espressione
nodale iniziale
\begin{equation*}
  V_S - V_A = V_A + sRCV_A + sRC(V_A - V_{BP})
\end{equation*}
sostituendo $V_A$
\begin{equation*}
  V_S + \frac{V_{BP}}{2sRC} = -\frac{V_{BP}}{2sRC} - sRC\frac{V_{BP}}{2sRC} - sRC\left(\frac{V_{BP}}{2sRC} + V_{BP}\right)
\end{equation*}
raccogliendo
\begin{equation*}
  V_S = -\left[\frac{1}{sRC} + 1 + sRC\right]V_{BP}
\end{equation*}
\begin{equation}
  \label{eq:notch-bp-sottrattore-tfunc-bp}
  \colorbox{blu}{\boxed{H_1(s) = -\frac{RCs}{1 + RCs + R^2C^2s^2}}}
\end{equation}

\noindent\colorbox{orange}{
  \begin{minipage}{1.0\linewidth}
    \paragraph{Richiesta (c)}
    Dalla restrizione di $A_{BP}$ sul dominio delle frequenze reali dedurre i valori attesi per i
    parametri del filtro (frequenza centrale, guadagno massimo e larghezza banda) e
    confrontarli con le misure di cui al punto precedente.
  \end{minipage}
}
\begin{equation*}
  \begin{cases}
    \mathcal{R}_{\omega}^N = 0 \\
    \mathcal{I}_{\omega}^N = RC\omega \\
    \mathcal{R}_{\omega}^D = 1 - R^2C^2\omega^2 \\
    \mathcal{I}_{\omega}^D = RC\omega
  \end{cases}
\end{equation*}
il guadagno
\begin{equation*}
  G_{BP}(\omega) = \frac{RC\omega}{\sqrt{(1 - R^2C^2\omega^2)^2 + (RC\omega)^2}}
\end{equation*}
per determinare il guadagno massimo
\begin{equation*}
  G_{BP}(\omega) = \frac{1}{\sqrt{1 + \left(\frac{1 - R^2C^2\omega^2}{RC\omega}\right)^2}}
\end{equation*}
guadagno massimo
\begin{equation}
  \label{eq:notch-bp-sottrattore-gain-max}
  \colorbox{blu}{\boxed{\max_{\omega \in [0, +\infty)} G_{BP}(\omega) = 1}}
\end{equation}
frequenza centrale
\begin{equation}
  \label{eq:notch-bp-sottrattore-frequenza-centrale}
  \colorbox{blu}{\boxed{f = \frac{1}{2\pi RC}}}
\end{equation}
frequenza di taglio
\begin{equation*}
  \begin{gathered}
    \frac{RC\omega}{\sqrt{(1 - R^2C^2\omega^2)^2 + (RC\omega)^2}} = \frac{1}{\sqrt2} \\ \implies
    (RC\omega)^2 = (1 - R^2C^2\omega^2)^2 \implies
    \begin{cases}
      +RC\omega_+ = 1 - R^2C^2\omega_+^2 \\
      -RC\omega_- = 1 - R^2C^2\omega_-^2
    \end{cases} \\ \implies
    \omega_+ = \frac{-1 \pm \sqrt{5}}{2RC} \quad \wedge \quad \omega_- = \frac{1 \pm \sqrt{5}}{2RC}
  \end{gathered}
\end{equation*}
una pulsazione ammissibile deve essere positiva, di conseguenza
\begin{equation*}
  \omega_{L, H} = \frac{\sqrt{5} \pm 1}{2RC}
\end{equation*}
attraverso la sostituzione $\omega \to 2\pi f$ si ottiene la larghezza di banda
\begin{equation}
  \label{eq:notch-bp-sottrattore-bw}
  \colorbox{blu}{\boxed{\Delta f = \frac{1}{2\pi RC}}}
\end{equation}

\noindent\colorbox{orange}{
  \begin{minipage}{1.0\linewidth}
    \paragraph{Richiesta (d)}
    Scrivere la funzione di trasferimento del blocco 2 (ovvero un'espressione di $V_{OUT}$ in
    termini di $V_S$ e $V_{BP}$) e ricavare infine la funzione di trasferimento complessiva del circuito
    ($A_{OUT}(s)=V_{OUT}(s)/V_S(s)$). Che tipo di circuito implementa il blocco 2?
  \end{minipage}
} \smallskip \\
Facendo uso del principio di sovvrapposizione per il quale
\begin{equation*}
  V_{OUT} = H_2(s)\Big|_{V_S = 0}V_{BP} + H_2(s)\Big|_{V_{BP} = 0}V_S
\end{equation*}
\begin{itemize}
\item Per $V_S = 0$
  \begin{equation*}
    \frac{V_{BP} - V_-}{R} = \frac{V_-}{R} + \frac{V_- - V_{OUT}}{R}
  \end{equation*}
  per il principio di corto-circuito virtuale dal fatto che $V_+ = 0$
  \begin{equation*}
    V_+ = V_- \implies \frac{V_{BP}}{R} = \frac{-V_{OUT}}{R}
  \end{equation*}
  da cui
  \begin{equation*}
    \boxed{H_2(s)\Big|_{V_S = 0} = -1}
  \end{equation*}
\item Per $V_{BP} = 0$
  \begin{equation*}
    \frac{-V_-}{R} = \frac{V_- - V_S}{R} + \frac{V_- - V_{OUT}}{R}
  \end{equation*}
  per il principio di corto-circuito virtuale dal fatto che $V_+ = 0$
  \begin{equation*}
    V_+ = V_- \implies 0 = \frac{-V_S}{R} + \frac{-V_{OUT}}{R}
  \end{equation*}
  da cui
  \begin{equation*}
    \boxed{H_2(s)\Big|_{V_{BP} = 0} = -1}
  \end{equation*}
\end{itemize}
infine
\begin{equation}
  \label{eq:notch-bp-sottrattore-tfunc-sum}
  \colorbox{blu}{\boxed{V_{OUT} = V_{BP} + V_S}}
\end{equation}
Il blocco 2 implementa un circuito sommatore. Inoltre siccome
$V_{BP}$ ha dipendenza da $V_S$ attraverso \ref{eq:notch-bp-sottrattore-tfunc-bp}
\begin{equation*}
  \begin{gathered}
    V_{OUT} = H_1(s)V_S + V_S = (1 + H_1(s))V_S \\ \implies
    H(s) \equiv 1 + H_1(s)
  \end{gathered}
\end{equation*}
la funzione di trasferimento complessiva
\begin{equation}
  \label{eq:notch-bp-sottrattore-tfunc}
  \colorbox{blu}{\boxed{H(s) = \frac{1 + R^2C^2s^2}{1 + RCs + R^2C^2s^2}}}
\end{equation}

\noindent\colorbox{orange}{
  \begin{minipage}{1.0\linewidth}
    \paragraph{Richiesta (e)}
    Calcolare i valori attesi del guadagno complessivo a frequenza $f_0$ ed ai limiti del dominio
    di frequenza ($f \to 0$ ed $f \to \infty$) e confrontarli con quanto misurato al punto 1f.
  \end{minipage}
} \smallskip \\
Il guadagno della funzione di trasferimento complessiva \ref{eq:notch-bp-sottrattore-tfunc}
\begin{equation*}
  G(\omega) = \frac{1 - R^2C^2\omega^2}{\sqrt{(1 - R^2C^2\omega^2)^2 + (RC\omega)^2}}
\end{equation*}
i valori attesi
\begin{itemize}
\item
  Per la frequenza $f_0 = 1/(2\pi RC)$ ottenuta da \ref{eq:notch-bp-sottrattore-frequenza-centrale},
  attraverso la sostituzione $\omega_0 \to f_0 / (2\pi)$, si ottiene il guadagno $\colorbox{blu}{\boxed{G(\omega_0) = G(1/RC) = 0}}$.
\item
  Per la frequenza $f \to 0$ il numeratore e denominatore si semplificano $\colorbox{blu}{\boxed{G(0) = 1}}$.
\item
  Per la frequenza $f \to \infty$ il termine dominante del numeratore \'e la potenza quadratica
  e sotto la radice il termine quartico, asintoticamente $\colorbox{blu}{\boxed{G(\omega \to \infty) = 1}}$.
\end{itemize}
Il risultato teorico \'e in accordo con i risultati precedenti, ora devono essere ripetuti i conti per la frequenza
centrale misurata.


\section*{Notch-T}

\noindent\colorbox{orange}{
  \begin{minipage}{1.0\linewidth}
    \paragraph{Richiesta (a)}
    Sulla base di semplici argomenti di natura qualitativa (ovvero senza effettuare il calcolo
    della funzione di trasferimento, come richiesto al punto successivo), mostrare che il
    circuito svolga effettivamente la funzione di notch, discutendone la risposta ai limiti del
    dominio delle frequenze reali.
  \end{minipage}
}
\begin{itemize}
\item
  Per $f \to 0$ i condensatori nel circuito ``aprono'' i loro rami. Di conseguenza il circuito
  finale si tratta di un op-amp in configurazione di \textit{inseguitore di tensione}, la tensione
  $V_S$ attraversa $R_1$ e $R_2$ in serie, nell'ipotesi di ``grossa'' impendenza in ingresso
  all'op-amp tali resistenze possono essere trascurate e nel regime lineare di operazione
  dell'op-amp vale
  \begin{equation*}
    V_{OUT} = A_d(V_S - V_{OUT}) \implies \colorbox{blu}{\boxed{V_{OUT} = \frac{A_d}{1+A_d}V_S}}
  \end{equation*}
\item
  Per $f \to \infty$ i condensatori $C_1$ e $C_2$ mettono ``in corto-circuito'' $V_S$ e $V_+$,
  nel regime lineare di operazione dell'op-amp vale
  \begin{equation*}
    V_{OUT} = A_d(V_S - V_{OUT}) \implies \colorbox{blu}{\boxed{V_{OUT} = \frac{A_d}{1+A_d}V_S}}
  \end{equation*}
\end{itemize}
La caratteristica del notch \'e quella dell'eliminazione di una banda all'uscita dal circuito,
data la topologia del circuito la condizione equivalente \'e quella dell'annullamento della
tensione all'ingresso \underline{non}-invertente $V_+$.
\begin{itemize}
\item
  \textit{Aumentando} la frequenza del segnale $f \to 0$, durante la caduta/salita di potenziale
  fra $R_1$ e $R_2$ il condensatore $C_3$ si scarica/carica, immettendo/assorbendo corrente da
  $R_2$ e di conseguenza \textit{ritardando} il segnale.
\item
  \textit{Diminuendo} la frequenza del segnale $f \to \infty$
  la carica del condensatore $C_1$ impiega tempo non nullo durante cui $C_2$ si carica pi\'u
  lentamente e di conseguenza \textit{ritardando} il segnale.
\end{itemize}
\noindent\colorbox{blu}{
  \begin{minipage}{1.0\linewidth}
    Siccome agli estremi del dominio di frequenza il segnale \'e senza attenuazione e per piccole
    variazioni di frequenza esso tenda ad attenuarsi \'e ``ragionevole'' assumere l'esistenza di un
    minimo anche detto \textit{notch}.
  \end{minipage}
}
\paragraph{Richiesta (b)}
Derivare analiticamente la funzione di trasferimento $A(s)=V_{OUT}(s)/V_S$ nell'ipotesi che
l'operazionale sia ideale ed assumendo per semplicit\'a che $R_1=R_2=R=10 kW$, $R3=R/2$,
$C_1=C_2=C=10 nF$ e $C_3=2C$. Si suggerisce di procedere:
\begin{enumerate}
\item[i.]
  considerando lo schema semplificato del circuito riportato in calce;
\item[ii.]
  ottenendo il suo equivalente per effetto della sostituzione di Th\'evenin dei blocchi A
  e B;
\item[iii.]
  quindi determinando la tensione nel nodo $V_+$ in termini di $V_S$ e $V_{OUT}$ ed in funzione
  della variabile complessa adimensionale $u = sRC$;
\item[iv.]
  infine, mettendo in relazione $V_+$ e $V_{OUT}$.
\end{enumerate}
\paragraph{Risposta (b.ii)}
\begin{itemize}
\item[a.]
  \begin{equation*}
    \begin{gathered}
      sC(V_S - V_A) = \frac{V_A - V_{OUT}}{R/2} \\
      sRC(V_S - V_A) = 2(V_A - V_{OUT}) \\
      sRCV_S + 2V_{OUT} = (2 + sRC)V_A \\ \implies
      V_{Th, A} \equiv V_A = \frac{sRCV_S + 2V_{OUT}}{2 + sRC}
    \end{gathered}
  \end{equation*}
  \begin{equation*}
    Z_{th, A} \equiv \frac{R}{2 + sRC}
  \end{equation*}
\item[b.]
  \begin{equation*}
    \begin{gathered}
      \frac{V_S - V_A}{R} = 2sCV_A \\
      V_{Th, B} \equiv V_A = \frac{V_S}{1 + 2sRC}
    \end{gathered}
  \end{equation*}
  \begin{equation*}
    Z_{Th, B} \equiv \frac{R}{1 + 2sRC}
  \end{equation*}
\end{itemize}
\paragraph{Risposta (b.iii)}
\begin{equation*}
  \begin{gathered}
    V_{Th, A} = \frac{uV_S + 2V_{OUT}}{2 + u} \qquad
    Z_{Th, A} = \frac{R}{2 + u} \\
    V_{Th, B} = \frac{V_S}{2 + u} \qquad
    Z_{Th, B} = \frac{R}{1 + 2u}
  \end{gathered}
\end{equation*}
\'E un partitore di tensione quindi
\begin{equation*}
  \begin{gathered}
    V_+ = \frac{Z_{Th, B} + R}{Z_{Th, A} + Z_{Th, B} + R + \frac{1}{sC}}V_{Th, A}
  \end{gathered}
\end{equation*}

\paragraph{Richiesta (c)}
Verificare che i limiti asintotici del guadagno (sia per $f \to 0$ che per $f \to \infty$) coincidano
con quanto osservato e qualitativamente dedotto al punto 2a.

\paragraph{Richiesta (d)}
Dall'espressione ottenuta dedurre i valori attesi per i parametri del filtro (frequenza
centrale, larghezza di banda e fattore di merito) e confrontarli con le misure di cui al
punto 1b (si consiglia di utilizzare la grandezza adimensionale $x = \omega RC = 2\pi fRC$
come variabile indipendente e di semplificare preliminarmente l'espressione della funzione
di trasferimento dividendo per il numeratore).

\paragraph{Richiesta (e)}
Osservare le propriet\'a di simmetria del plot di Bode e giustificarle alla luce
del fatto che $A(1/x) = A^*(x)$.


\section*{Filtro Bandpass (retroazione multipla)}

\noindent\colorbox{orange}{
  \begin{minipage}{1.0\linewidth}
    \paragraph{Richiesta (a)}
    Dimostrare, sulla base di semplici argomenti di natura qualitativa, che il circuito proposto
    svolga effettivamente la funzione di filtro passa-banda, discutendone la risposta ai limiti
    del dominio delle frequenze reali.
  \end{minipage}
}
\begin{itemize}
\item
  Per $f \to 0$ i condensatori $C_1$ e $C_2$ aprono i loro rami, le tensioni d'ingresso all'op-amp
  sono nulle
  \begin{equation*}
    V_+ = 0, ~V_- = 0 \implies \colorbox{blu}{\boxed{V_{OUT} = 0}}
  \end{equation*}
\item
  Per $f \to \infty$ il condensatore $C_2$ mette in ``corto-circuito'' $V_A$ con $V_{OUT}$
  \begin{equation*}
    V_A = V_{OUT}
  \end{equation*}
  il condensatore $C_1$ mette in ``corto-circuito'' $V_A$ con $V_-$
  \begin{equation*}
    V_{OUT} = A_d(0-V_A) = A_d(0-V_{OUT}) \implies \colorbox{blu}{\boxed{V_{OUT} = \frac{0}{1 + A_d} = 0}}
  \end{equation*}
\end{itemize}

\noindent\colorbox{orange}{
  \begin{minipage}{1.0\linewidth}
    \paragraph{Richiesta (b)}
    Derivare analiticamente la funzione di trasferimento complessa $A_v(s)=V_{OUT}(s)/V_S(s)$
    nell’ipotesi che l’operazionale sia ideale (\textit{suggerimento: si determini la tensione nel
      punto A del circuito in termini di $V_S$ e $V_{OUT}$; quindi si esprima $V_A$ in termini di $V_{OUT}$
      assumendo che l'op-amp funzioni in regime lineare ed utilizzando il principio di cortocircuito
      virtuale; si eguaglino infine le espressioni ottenute per ottenere la funzione di trasferimento
    }).
    Per semplicità di calcolo, si assuma inoltre (ipotesi compatibile con i
    valori e le incertezze nominali dei componenti) che $R_1=R_2=R$, $R_3=100R$, $C_1=C_2=C$.
  \end{minipage}
}

\paragraph{Risposta (b)}
Seguendo il suggerimento della richiesta, attraverso la \textit{KCL}
si ottiene la seguente relazione per la tensione $V_A$ in termini di $V_S$
e $V_{OUT}$
\begin{equation*}
  \frac{V_S - V_A}{R} = sC(V_A - V_{OUT}) + sC(V_A - V_-) + \frac{V_A}{R}
\end{equation*}
attraverso la \textit{KCL} si ottiene anche una seconda relazione per $V_A$
in termini di $V_{OUT}$
\begin{equation*}
  sC(V_A - V_-) = \frac{V_- - V_{OUT}}{100R}
\end{equation*}
siccome $V_+ = 0$ nell'ipotesi di validit\'a per il principio di corto-cicuito
virtuale vale
\begin{equation*}
  V_+ = V_- \implies sCV_A = -\frac{V_{OUT}}{100R}
\end{equation*}
la soluzione del sistema dell'equazione precedente appaiata con la prima espressione
ricavata coinvolge \textit{unicamente} $V_S$ e $V_{OUT}$. Il primo passaggio
\'e una manipolazione algebrica della prima espressione e la sostituzione di $V_- = 0$
\begin{equation*}
  V_S - V_A = sRC(V_A - V_{OUT}) + sRCV_A + V_A
\end{equation*}
si raccolgono a destra i termini con $V_A$
\begin{equation*}
  V_S = 2(1 + sRC)V_A - sRCV_{OUT}
\end{equation*}
si sostituisce $V_A$ e raccoglie nuovamente
\begin{equation*}
  V_S = -\frac{1}{50sRC}\left(1 + sRC + 50s^2R^2C^2\right)V_{OUT}
\end{equation*}
finalmente si ottiene la funzione di trasferimento
\begin{equation}
  \label{eq:bp-multi-feedback-tfunc}
  \colorbox{blu}{\boxed{H(s) = -\frac{50sRC}{1 + sRC + 50s^2R^2C^2}}}
\end{equation}

\noindent\colorbox{orange}{
  \begin{minipage}{1.0\linewidth}
    \paragraph{Richiesta (c)}
    Dalla funzione di trasferimento dedurre i valori attesi per i parametri del filtro (frequenza
    centrale, guadagno massimo e larghezza banda) e confrontarli con le misure di cui ai
    punti 1b, 1c, 1d.
  \end{minipage}
}
\paragraph{Risposta (c)}
La funzione di trasferimento \ref{eq:bp-multi-feedback-tfunc} \'e in forma canonica, calcolando
le parti reali ed immaginarie del numeratore e denominatore con dominio ristretto all'asse
immaginario si ottiene
\begin{equation*}
  \begin{cases}
    \mathcal{R}_{\omega}^N = 0 \\
    \mathcal{I}_{\omega}^N = -50RC\omega \\
    \mathcal{R}_{\omega}^D = 1 - 50R^2C^2\omega^2 \\
    \mathcal{I}_{\omega}^D = RC\omega
  \end{cases}
\end{equation*}
di conseguenza ottengo la funzione di guadagno mediante \ref{eq:da-tfunc-a-gain}
\begin{equation*}
  G(\omega) = \frac{50RC\omega}{\sqrt{(1 - 50R^2C^2\omega^2)^2 + (RC\omega)^2}}
\end{equation*}
dividendo numeratore e denominatore per $RC\omega$ si ottiene un espressione pi\'u
semplice da cui ricavare il guadagno massimo
\begin{equation*}
  G(\omega) = \frac{50}{\sqrt{1 + \left(\frac{1 - 50R^2C^2\omega^2}{RC\omega}\right)^2}}
\end{equation*}
Il valore minimo che l'espressione sotto la radice potrebbe idealmente assumere \'e $1$,
questo valore \'e ottenibile per $1 - 50R^2C^2\omega^2 = 0$ dunque
\begin{equation}
  \label{eq:bp-multi-feedback-gain-max}
  \colorbox{blu}{\boxed{\max_{\omega \in [0, \infty)}G(\omega) = \frac{50}{\sqrt{1 + 0}} = 50}}
\end{equation}
inoltre attraverso la sostituzione $\omega \to f = 2\pi \omega$ si ottiene la frequenza centrale
\begin{equation}
  \label{eq:bp-multi-feedback-frequenza-centrale}
  \colorbox{blu}{\boxed{f = \frac{1}{10\pi\sqrt{2}RC} \approx 2250~\text{Hz}}}
\end{equation}

\noindent\colorbox{orange}{
  \begin{minipage}{1.0\linewidth}
    \paragraph{Richiesta (d)}
    Dare un’interpretazione qualitativa (e se possibile quantitativa) delle caratteristiche del
    segnale di uscita osservato al punto 1e. Che tipo di segnale appare? A quale ampiezza?
  \end{minipage}
}

\paragraph{Risposta (d)}
Se si invia nell'ingresso un onda quadra di frequenza pari alla frequenza di centro-banda, siccome
l'onda quadra \'e rappresentabile come una combinazione lineare di onde sinusoidali la cui prima
componente a frequenza non-nulla \'e a frequenza di centro-banda. Il risultato dovrebbe essere
un unica onda sinusoidale
\begin{equation}
  \colorbox{blu}{\boxed{V_{OUT} = 50\cdot\frac{4\cdot 50~\text{mV}}{\pi} \approx 3.18~\text{V}}}
\end{equation}



\section*{Filtro Lowpass (retroazione multipla)}

\noindent\colorbox{orange}{
  \begin{minipage}{1.0\linewidth}
    \paragraph{Richiesta (a)}
    Dimostrare, sulla base di semplici argomenti di natura qualitativa, che il circuito proposto
    svolga effettivamente la funzione di filtro passa-basso, discutendone la risposta ai limiti
    del dominio delle frequenze reali.
  \end{minipage}
}
\paragraph{Risposta (a)}
\begin{itemize}
\item
  Per $f \to 0$ i rami di $C_1$ e $C_2$ si aprono, $R_2$ \'e dunque in serie con l'impedenza
  d'ingresso ``grande'' dell'op-amp, \'e trascurabile. Il circuito assume la forma di un
  amplificatore invertente.
\item
  Per $f \to +\infty$ il condensatore $C_2$ ``corto-circuita'' $V_+$ con $V_{OUT}$ per cui
  \begin{equation*}
    V_{OUT} = A_d(0 - V_{OUT}) \implies \colorbox{blu}{\boxed{V_{OUT} = \frac{0}{1 + A_bd} = 0}}
  \end{equation*}
\end{itemize}

\noindent\colorbox{orange}{
  \begin{minipage}{1.0\linewidth}
    \paragraph{Richiesta (b)}
    Derivare analiticamente la funzione di trasferimento complessa $A_v(s)=V_{OUT}(s)/V_S(s)$
    nell'ipotesi che l'operazionale sia ideale
    (\textit{suggerimento: si determini la tensione nel nodo A del circuito in termini di $V_S$ e
      $V_{OUT}$; quindi si esprima $V_A$ in termini di $V_{OUT}$ assumendo che l'op-amp funzioni in
      regime lineare ed utilizzando il principio di cortocircuito virtuale; si eguaglino
      infine le espressioni ottenute per ottenere la funzione di trasferimento}).
    Per semplicit\'a di calcolo, si assuma inoltre (ipotesi compatibile con i
    valori e le incertezze nominali dei componenti) che $R_1=R_3=R$, $R_2=R/2$, $C_2=C$, $C_1=2C$.
  \end{minipage}
}
\paragraph{Risposta (b)}
Usando la \textit{KCL} al nodo A
\begin{equation*}
  \frac{V_S - V_A}{R} = \frac{V_A - V_{OUT}}{R} + \frac{V_A - V_-}{R/2} + 2sCV_A
\end{equation*}
Usando la \textit{KCL} al nodo + e l'ipotesi di corto-circuito virtuale sapendo $V_+ = 0$
\begin{equation*}
  \begin{gathered}
    \frac{V_A - V_-}{R/2} = sC(V_- - V_{OUT}) \\ \overset{V_+ = V_-}{\implies}
    \frac{V_A}{R/2} = -sCV_{OUT}
  \end{gathered}
\end{equation*}
applicando il principio di corto-circuito virtuale alla prima espressione
\begin{equation*}
  \frac{V_S - V_A}{R} = \frac{V_A - V_{OUT}}{R} + \frac{V_A}{R/2} + 2sCV_A
\end{equation*}
moltiplicando LHS e RHS per $R$
\begin{equation*}
  V_S - V_A = V_A - V_{OUT} + 2V_A + 2sRCV_A
\end{equation*}
raccogliendo i termini $V_A$ a destra
\begin{equation*}
  V_S = 2(2 + sRC)V_A - V_{OUT}
\end{equation*}
sostituisco $V_A$
\begin{equation*}
  V_S = -(1 + 2sRC + s^2R^2C^2)V_{OUT}
\end{equation*}
\begin{equation}
  \label{eq:lp-multi-feedback-tfunc}
  \colorbox{blu}{\boxed{H(s) = -\frac{1}{(1 + sRC)^2}}}
\end{equation}

\noindent\colorbox{orange}{
  \begin{minipage}{1.0\linewidth}
    \paragraph{Richiesta (c)}
    Dalla funzione di trasferimento dedurre le caratteristiche del filtro osservate al punto 1)
    (guadagno a centro-banda, frequenza naturale, andamento di modulo/fase in funzione
    della frequenza) e confrontarle con quanto misurato ai punti 1c, 1d, 1e.
  \end{minipage}
}
\paragraph{Risposta (c)}
La funzione di trasferimento \ref{eq:lp-multi-feedback-tfunc} \'e in forma canonica, calcolando
le parti reali ed immaginarie del numeratore e denominatore con dominio ristretto all'asse
immaginario si ottiene
\begin{equation*}
  \begin{cases}
    \mathcal{R}_{\omega}^N = -1 \\
    \mathcal{I}_{\omega}^N = 0 \\
    \mathcal{R}_{\omega}^D = 1 - R^2C^2\omega^2 \\
    \mathcal{I}_{\omega}^D = 2RC\omega
  \end{cases}
\end{equation*}
di conseguenza ottengo la funzione di guadagno mediante \ref{eq:da-tfunc-a-gain}
\begin{equation*}
  G(\omega) = \frac{1}{\sqrt{(1 - R^2C^2\omega^2)^2 + (2RC\omega)^2}} = \frac{1}{1 + R^2C^2\omega^2}
\end{equation*}
l'andamento del guadagno in termini della frequenza si ottiene sostituendo $\omega \to 2\pi f$
\begin{equation}
  \label{eq:lp-multi-feedback-gain}
  \colorbox{blu}{\boxed{G(f) = \frac{1}{1 + 4\pi^2R^2C^2f^2}}}
\end{equation}
il guadagno massimo \'e chiaramente per $\omega \to 0$
\begin{equation}
  \label{eq:lp-multi-feedback-gain-max}
  \colorbox{blu}{\boxed{\max_{\omega \in [0, +\infty)}G(\omega) = 1}}
\end{equation}
di conseguenza la frequenza centrale
\begin{equation}
  \label{eq:lp-multi-feedback-frequenza-centrale}
  \colorbox{blu}{\boxed{f = 0}}
\end{equation}
la frequenza di taglio si ottiene
\begin{equation*}
  \frac{1}{1 + R^2C^2\omega^2} = \frac{1}{\sqrt{2}} \implies
  \omega = \pm\frac{\sqrt{\sqrt{2} - 1}}{RC}
\end{equation*}
la pulsazione e frequenza ammettono solamente valori positivi, sostituendo $\omega \to 2\pi f$
\begin{equation}
  \label{eq:lp-multi-feedback-frequenza-cut}
  \colorbox{blu}{\boxed{f = \frac{\sqrt{\sqrt{2} - 1}}{2\pi RC}}}
\end{equation}
la frequenza di taglio \'e diversa dalla frequenza naturale
\begin{equation}
  \label{eq:lp-multi-feedback-frequenza-naturale}
  \colorbox{blu}{\boxed{f = \frac{1}{2\pi RC}}}
\end{equation}
l'andamento della fase richiede identificare il corretto semi-piano, il segno della parte reale
di $H(s)$
\begin{equation*}
  R^2C^2\omega^2 - 1
\end{equation*}
non \'e ben determinato nell'intervall $\omega \in [0, \infty)$, mentre
\begin{equation*}
  2RC\omega > 0
\end{equation*}
\'e ben determinato nell'intervall $\omega \in [0, \infty)$, l'andamento della
fase (in funzione della pulsazione \'e dato da) \ref{eq:da-tfunc-a-fase-immaginario-positivo}
attraverso la sostituzione $\omega \to 2 \pi f$ si ottiene
\begin{equation}
  \label{eq:lp-multi-feedback-fase}
  \colorbox{blu}{\boxed{\phi(f) = \frac{\pi}{2} + \arctan\left(\frac{1 - 4\pi^2R^2C^2f^2}{4\pi RC f}\right)}}
\end{equation}

\noindent\colorbox{orange}{
  \begin{minipage}{1.0\linewidth}
    \paragraph{Richiesta (d)}
    Dare un'interpretazione qualitativa (e se possibile quantitativa) delle caratteristiche del
    segnale di uscita osservato al punto 1f. Che tipo di segnale appare? A quale ampiezza?
  \end{minipage}
}
\paragraph{Risposta (d)}
Se si invia nell'ingresso un onda quadra di frequenza pari alla frequenza di centro-banda, siccome
l'onda quadra \'e rappresentabile come una combinazione lineare di onde sinusoidali la cui prima
componente a frequenza non-nulla \'e a frequenza di centro-banda. Il risultato dovrebbe essere
un unica onda sinusoidale
\begin{equation}
  \colorbox{blu}{\boxed{V_{OUT} = G(5~\text{kHz})\cdot\frac{4}{\pi}\cdot 5~\text{V}\approx 586~\text{mV}}}
\end{equation}
il risultato dovrebbe essere in accordo con quanto misurato sperimentalmente.


\section{Oscillatore Fase 0}

\noindent\colorbox{orange}{
  \begin{minipage}{1.0\linewidth}
    \paragraph{Richiesta (a)}
    Discutere le caratteristiche del blocco 1 e determinare il suo guadagno $A_1$, al momento
    trascurando la presenza dei diodi.
  \end{minipage}
}
\paragraph{Risposta (a)}
Il ruolo del blocco $A_1$ (ignorando la presenza di diodi) \'e un comune amplificatore a schema
\underline{non}-invertente con op-amp, la funzione di trasferimento \'e nota.
\begin{equation}
  \label{eq:oscillatore-fase-0-tfunc-1}
  \colorbox{blu}{\boxed{H_1(s) = \frac{R + (R/5 + R)}{R} = \frac{11}{5}}}
\end{equation}

\noindent\colorbox{orange}{
  \begin{minipage}{1.0\linewidth}
    \paragraph{Richiesta (b)}
    Determinare la funzione di trasferimento $A_2(s)$ del blocco 2.
  \end{minipage}
}
\paragraph{Risposta (b)}
Il blocco 2 \'e composto da un \textit{integratore passivo} in serie con un \textit{inseguitore di tensione}
con op-amp ed un \textit{derivatore passivo}.
\begin{equation}
  \label{eq:oscillatore-fase-0-tfunc-2}
  \colorbox{blu}{\boxed{H_2(s) = \frac{sRC}{(1 + sRC)^2}}}
\end{equation}

\noindent\colorbox{orange}{
  \begin{minipage}{1.0\linewidth}
    \paragraph{Richiesta (c)}
    Verificare che esiste un unico valore di frequenza per cui il loop-gain $A_1A_2$ soddisfi la
    condizione di Barkhausen e metterlo in relazione alla frequenza di oscillazione
    precedentemente misurata.
  \end{minipage}
}
\paragraph{Risposta (c)}
Il loop gain si ottiene
\begin{equation*}
  L(s) \equiv H_1(s)H_2(s) = \frac{11}{5}\frac{sRC}{(1 + sRC)^2}
\end{equation*}
nello studio della fase il fattore positivo non contribuisce all'identificazione del semi-piano
di operazione dunque sar\'a sufficiente studiare $H_2(s)$.
Il segno del semi-piano reale
\begin{equation*}
  2R^2C^2\omega^2 > 0
\end{equation*}
il segno del semi-piano complesso
\begin{equation*}
  RC\omega(1 - R^2C^2\omega^2)
\end{equation*}
non \'e ben definito. Di conseguenza si usa \ref{eq:da-tfunc-a-fase-reale-positivo}
\begin{equation}
  \label{eq:oscillatore-fase-0-loop-gain-fase}
  \colorbox{blu}{\boxed{\phi(\omega) = \arctan\left(\frac{1 - R^2C^2\omega^2}{2RC\omega}\right)}}
\end{equation}
la frequenza \'e direttamente proporzionale alla pulsazione, la fase si annulla quando
\begin{equation*}
  1 - R^2C^2\omega^2 = 0
\end{equation*}
siccome le pulsazione ammesse sono positive esiste un unica frequenza che soddisfa la
condizione di barkhausen, la frequenza centrale (ottenuta sostituendo $\omega \to 2\pi f$)
\begin{equation}
  \label{eq:oscillatore-fase-0-loop-gain-frequenza-centrale}
  \colorbox{blu}{\boxed{f = \frac{1}{2\pi RC}}}
\end{equation}

\noindent\colorbox{orange}{
  \begin{minipage}{1.0\linewidth}
    \paragraph{Richiesta (d)}
    Discutere qualitativamente e quantitativamente il ruolo dei diodi alla luce dell’andamento
    osservato (punti 1c-e) del loop-gain in funzione dell’ampiezza di $V_1$, in particolare
    determinando un valore atteso per l’ampiezza dell’aut
  \end{minipage}
}
\paragraph{Risposta (d)}
La funzione di trasferimento del blocco 1 modificata dall'operazione dei diodi, dove
$Z = D||R||D$.
\begin{equation*}
  H_1(s, Z) = \frac{R + (R/5 + Z)}{R} = \frac65 + \frac{Z}{R}
\end{equation*}
di conseguenza la funzione di loop-gain modificata
\begin{equation*}
  L(s, Z) = \left(\frac65 + \frac{Z}{R}\right)\frac{sRC}{(1 + sRC)^2}
\end{equation*}
la legge sulla fase non viene modificata siccome il termine modificato \'e sempre
un fattore di $H_2(s)$, la pulsazione/frequenza centrale rimane invariata, occorre studiare
il guadagno che per soddisfare la condizione di Barkhausen deve essere unitario, utilizzando
\ref{eq:da-tfunc-a-gain}
\begin{equation}
  \label{eq:oscillatore-fase-0-resistenza-dinamica}
  \begin{gathered}
    G(\omega, Z) = \left(\frac65 + \frac{Z}{R}\right)\frac{RC\omega}{\sqrt{(1 - R^2C^2\omega^2)^2 + (2RC\omega)^2}} = 1 \\ \implies
    \left(\frac65 + \frac{Z}{R}\right)\frac{1}{\sqrt{(0)^2 + (2)^2}} = 1 \implies \colorbox{blu}{\boxed{Z = \frac45 R}}
  \end{gathered}
\end{equation}

\noindent\colorbox{orange}{
  \begin{minipage}{1.0\linewidth}
    \paragraph{Richiesta (e)}
    Enfatizzare il ruolo svolto da $U_2$ discutendo come sarebbe cambiata la risposta alla domanda
    3b eliminandolo (ovvero corto-circuitando i terminali delle capacità) e come sarebbe stato
    necessario in questo caso modificare il blocco 1 per ripristinare la condizione di Barkhausen.
  \end{minipage}
}
\paragraph{Risposta (e)}
Rimuovendo l'op-amp $U_2$ l'integratore passivo e il derivatore passivo sono sempre in serie ma sono
accoppiati, non la tensione intermedia non \'e pi\'u proveniente da un partitore di tensione
\begin{equation*}
  \begin{gathered}
    \frac{V_S - V}{R} = sCV + sC(V - V_{OUT}) \\
    V_{OUT} = \frac{sRC}{1 + sRC}V
  \end{gathered}
\end{equation*}
isolando il termine $V_S$ a sinistra
\begin{equation*}
  \begin{gathered}
    V_S = (1 + 2sRC)V - sRCV_{OUT}
  \end{gathered}
\end{equation*}
sostituendo $V \to \frac{1 + sRC}{sRC}V_{OUT}$ e $sRCV \to (1+sRC)V_{OUT}$
\begin{equation*}
  \begin{gathered}
    V_S = \left[\frac{1 + sRC}{sRC} + 2(1 + sRC) - sRC\right]V_{OUT}
  \end{gathered}
\end{equation*}
moltiplicando entrambi i lati per $sRC$ si ottiene
\begin{equation}
  \label{eq:oscillatore-fase-0-tfunc-2-mod}
  \colorbox{blu}{\boxed{\tilde{H}_2(s) = \frac{sRC}{1 + 3sRC + s^2R^2C^2}}}
\end{equation}
Ripetendo i conti per il calcolo della pulsazione a fase nulla ci si rende conto
che essa rimane invariata (nell'espressione all'interno dell'$\arctan$ cambia solo
che a denominatore $2RC\omega \to 3RC\omega$).
Si modifica per\'o il guadagno e ripartendo da \ref{eq:oscillatore-fase-0-resistenza-dinamica}
si ottiene la nuova condizione sull'operazione dei diodi perch\'e sia soddisfatta la
condizione di Barkhausen
\begin{equation}
  \label{eq:oscillatore-fase-0-resistenza-dinamica-mod}
  \begin{gathered}
    \tilde{G}(\omega, \tilde{Z}) = \left(\frac65 + \frac{\tilde{Z}}{R}\right)\frac{RC\omega}{\sqrt{(1 - R^2C^2\omega^2)^2 + (3RC\omega)^2}} = 1 \\ \implies
    \left(\frac65 + \frac{\tilde{Z}}{R}\right)\frac{1}{\sqrt{(0)^2 + (3)^2}} = 1 \implies \colorbox{blu}{\boxed{\tilde{Z} = \frac95 R}}
  \end{gathered}
\end{equation}


\section*{Oscillatore integratore-derivatore}
\noindent\colorbox{orange}{
  \begin{minipage}{1.0\linewidth}
    \paragraph{Richiesta (a)}
    Discutere, sulla base di semplici argomenti di natura qualitativa (ovvero senza effettuare
    i calcoli richiesti ai punti successivi), la funzione svolta dai blocchi 1 (ingresso $V_1$,
    uscita $V_2$) e 2 (ingresso $V_2$, uscita $V_3$) del circuito.
  \end{minipage}
}
\paragraph{Risposta (a)}
\begin{itemize}
\item
  Nel blocco 1 \'e immediato notare che per $f \to 0$ il condensatore $C_1$ apre il ramo
  e la tensione d'uscita $V_2$ \'e nulla, mentre per $f \to \infty$ il condensatore $C_1$
  ``corto-circuita'' il ramo e il circuito assume lo schema di \textit{amplificatore invertente};
  il blocco 1 \'e un circuito a schema di \colorbox{blu}{\textit{derivatore reale con op-amp}}.
\item
  Nel blocco 2 $f \to 0$ il condensatore $C_1$ apre il ramo e il circuito assume lo schema di
  \textit{amplificatore invertente}, mentre per $f \to \infty$
  \begin{equation*}
    V_{OUT} = A_d(0 - V_{OUT}) \implies V_{OUT} = \frac{0}{1 + A_d} = 0
  \end{equation*}
  il blocco 2 \'e un circuito a schema di \colorbox{blu}{\textit{integratore reale con op-amp}}.
\end{itemize}

\noindent\colorbox{orange}{
  \begin{minipage}{1.0\linewidth}
    \paragraph{Richiesta (b)}
    Scrivere un'espressione analitica per le funzioni di trasferimento $A_1$, $A_2$ dei singoli
    blocchi e del loop-gain, nelle ipotesi di idealit\'a dell'operazionale e (al momento) di
    polarizzazione inversa dei diodi, e giustificare perché quest'ultima possa essere ottenuta
    come prodotto delle precedenti.
  \end{minipage}
} \smallskip \\
Nell'ipotesi che nessuna corrente scorra fra i diodi, la ramo di feedback dell'op-amp U1 ha
impedenza equivalente $\tilde{R}$ dove
\begin{equation*}
  \tilde{R} \approx 2R + R/2 = \frac52 R
\end{equation*}
lo schema del blocco $A_1$ \'e ben noto, la funzione di trasferimento
\begin{equation}
  \label{eq:oscillatore-integratore-derivatore-tfunc-1}
  \colorbox{blu}{\boxed{H_1(s) = -\frac{s\tilde{R}C}{1 + sRC}}}
\end{equation}
lo schema del blocco $A_2$ \'e ben noto, la funzione di trasferimento
\begin{equation}
  \label{eq:oscillatore-integratore-derivatore-tfunc-2}
  \colorbox{blu}{\boxed{H_2(s) = -\frac{1}{1 + sRC}}}
\end{equation}
\begin{equation}
  \label{eq:oscillatore-integratore-derivatore-loop-gain}
  \colorbox{blu}{\boxed{L(s) \equiv H_1(s)H_2(s) = \frac{s\tilde{R}C}{(1 + sRC)^2}}}
\end{equation}
la funzione di trasferimento finale \'e compatibile con la funzione
di trasferimento di un \textit{filtro passa-banda}. \\

\noindent\colorbox{orange}{
  \begin{minipage}{1.0\linewidth}
    \paragraph{Richiesta (c)}
    Mostrare che esiste unica una frequenza tale che il loop-gain, calcolato sul dominio delle
    frequenze reali, abbia fase nulla e confrontare la compatibilit\'a di detta frequenza con
    quelle misurate ai punti 1a ed 1c.
  \end{minipage}
}
\paragraph{Risposta (c)}
Lo studio della fase di $L(s)$ richiede di individuare il semi-piano sui cui \'e definita la parte
reale ed immaginaria della fase
\begin{equation*}
  \begin{cases}
    \mathcal{R}_{\omega}^N = 0 \\
    \mathcal{I}_{\omega}^N = \tilde{R}C\omega \\
    \mathcal{R}_{\omega}^D = 1 - R^2C^2\omega^2 \\
    \mathcal{I}_{\omega}^D = 2RC\omega
  \end{cases}
\end{equation*}
si studia il segno del semi-piano reale di $L(s)$
\begin{equation*}
  2R\tilde{R}C\omega > 0
\end{equation*}
si studia il segno del semi-piano immaginario di $L(s)$
\begin{equation*}
  \tilde{R}C\omega(1 - R^2C^2\omega^2)
\end{equation*}
il segno non \'e definito, di conseguenza la legge di fase per $L(s)$ \'e definita da
\ref{eq:da-tfunc-a-fase-reale-positivo} e la si fa dipendere dalla frequenza con la sostituzione
$\omega \to 2\pi f$
\begin{equation}
  \label{eq:oscillatore-integratore-derivatore-fase}
  \colorbox{blu}{\boxed{\phi(f) = \arctan\left(\frac{1 - 4\pi^2R^2C^2f^2}{4\pi RCf}\right)}}
\end{equation}
(il fattore moltiplicativo che differisce $\tilde{R}$ da $R$ si vede che non contribuisce).
L'annullamento della fase richiede l'annullamento del numeratore
\begin{equation}
  \label{eq:oscillatore-integratore-derivatore-frequenza-centrale}
  \colorbox{blu}{\boxed{f = \frac{1}{2\pi RC} \approx 1592~\text{Hz}}}
\end{equation}
essendo la frequenza una quantit\'a positiva la soluzione \'e \textit{unica}.
La stima \'e qualitativamente compatibile con quanto osservato dal circuito simulato tramite
\textit{falstad}. \\

\noindent\colorbox{orange}{
  \begin{minipage}{1.0\linewidth}
    \paragraph{Richiesta (d)}
    Calcolare la resistenza equivalente del parallelo $R_4$-$D_1$-$D_2$ tale che a quella frequenza si
    realizzi la condizione di Barkhausen.
  \end{minipage}
}
\paragraph{Risposta (d)}
Modificando la risposta di \ref{eq:oscillatore-integratore-derivatore-tfunc-1} tenendo in considerazione la
rete di diodi si ottiene
\begin{equation*}
  \begin{gathered}
    Z \equiv D_1||R||D_2 \\
    H_1(s) = -\frac{s(\frac32 R + Z)C}{1 + sRC} \\ \implies
    L(s) \equiv H_1(s)H_2(s) = \frac{s(\frac32 R + Z)C}{(1 + sRC)^2}
  \end{gathered}
\end{equation*}
per quanto osservato prima la frequenza centrale non dipende da $\tilde{R}$ di conseguenza bisogner\'a
studiare la \textit{condizione di barkhausen} sul modulo, utilizzando \ref{eq:da-tfunc-a-gain} si ottiene
\begin{equation*}
  G(\omega) = \frac{(\frac32 R + Z)C\omega}{\sqrt{(1 - R^2C^2\omega^2)^2 + (2RC\omega)^2}} = 1
\end{equation*}
si impone la pulsazione in fase usando \ref{eq:oscillatore-integratore-derivatore-frequenza-centrale}
e sostituendo $f \to \frac{\omega}{2\pi}$
\begin{equation*}
  \begin{gathered}
    \frac{(\frac32 R + Z)\frac{1}{R}}{\sqrt{0^2 + (2)^2}} = 1 \implies
    \colorbox{blu}{\boxed{Z = \frac12 R}}
  \end{gathered}
\end{equation*}

\noindent\colorbox{orange}{
  \begin{minipage}{1.0\linewidth}
    \paragraph{Richiesta (e)}
    Sulla base della misura nelle condizioni di cui al punto 1e provare a determinare una
    stima per il valore atteso dell'ampiezza del segnale di oscillazione in $V_2$.
  \end{minipage}
}
\paragraph{Risposta (e)}
Sulla base del risultato precedente, la funzione di trasferimento del primo blocco vale
\begin{equation*}
  H_1(s) = -\frac{s(\frac32 R + \frac12 R)C}{1 + sRC} = -\frac{2sRC}{1 + sRC}
\end{equation*}
il guadagno si ricava usando \ref{eq:da-tfunc-a-gain} per cui
\begin{equation*}
  G_1(\omega) = \frac{2RC\omega}{\sqrt{1 + (RC\omega)^2}}
\end{equation*}
il guadagno a centro banda
\begin{equation*}
  G_1(1/(RC)) = \frac{2}{\sqrt{1 + 1^2}} = \sqrt2
\end{equation*}
per una tensione $V_3 = 1.377~\text{V}$ il valore di aspettazione di $V_2$
\begin{equation*}
  \colorbox{blu}{\boxed{V_2 = V_3\sqrt{2} \approx 1.947~\text{V}}}
\end{equation*}



% Si assumano tutti gli operazionali ideali ed operanti in regime lineare. Si considerino i valori
delle resistenze pari a quelli nominali.

\paragraph{Richiesta (a)}
Mostrare sulla base di semplici argomenti che i diodi D1 e D2 non possono essere
simultaneamente nello stesso stato di polarizzazione (se cioè uno è in conduzione l’altro deve
essere interdetto e viceversa).
\paragraph{Risposta (a)}
\begin{quote}
  ``Entrambi i diodi sono entrambi in conduzione?''
\end{quote}
Assumendo senza perdita di generalit\'a una tensione d'ingresso \textit{positiva}
sul canale $V_S$ allora la tensione fra le resistenze e i diodi dovranno essere di segno
\textit{concorde} siccome gli elementi resistivi non possiedono reattanza. Assumendo per
assurdo che entrambi i diodi siano in conduzione si ottiene uno schema di amplificatore
\textit{invertente}, dunque la tensione all'uscita dell'amplificatore $V_{out}$ avr\'a
segno \textit{discorde} rispetto tensione d'ingresso $V_S$. Il diodo ``al di sotto''
dell'ingresso non-invertente non pu\'o essere messo in conduzione perch\'e all'ingresso
avr\'a potenziale \textit{negativo} e all'uscita potenziale \textit{positivo}, una
contraddizione (un argomento analogo viene fatto assumendo la tensione d'ingresso \textit{negativa}).

\begin{quote}
  ``Entrambi i diodi sono entrambi interdetti?''
\end{quote}
Se si assumesse invece che nessuno dei diodi \'e in conduzione allora la tensione all'uscita
dell'amplificatore \'e nulla, quindi nuovamente almeno uno fra i diodi entra in conduzione
(perch\'e la tensione d'ingresso \'e non nulla), un'altra contraddizione.
\textit{Almeno uno} dei diodi dovr\'a essere in conduzione.

\paragraph{Richiesta (b)}
Alla luce del precedente argomento si giustifichino le identità tra le semi-onde di VA, VB, VS
di cui al punto 1b.
\paragraph{Risposta (b)}
Siccome un diodo alla volta pu\'o entrare in conduzione le possibili tensioni all'uscite
di $V_A$ e $V_B$ sono rispettivamente
\begin{equation*}
  \begin{gathered}
    V_S > 0 \implies \frac{V_S - V_{-U1}}{R_1} = \frac{V_{-U1} - V_{+U3}}{R_2} \\
    V_S < 0 \implies \frac{V_S - V_{-U1}}{R_1} = \frac{V_{-U1} - V_{+U4}}{R_3}
  \end{gathered}
\end{equation*}
siccome $R_1 = R_2 = R_3$ e $V_{-U1} = V_{+U1} = 0$ allora
\begin{equation*}
  \begin{gathered}
    V_S > 0 \implies
    \begin{cases}
      V_{+U3} = -V_S \\
      V_{+U4} = 0
    \end{cases} \\
    V_S < 0 \implies
    \begin{cases}
      V_{+U3} = 0 \\
      V_{+U4} = -V_S
    \end{cases}
  \end{gathered} \implies
  V_{+U3} + V_{+U4} = -V_S
\end{equation*}
inoltre $V_{+U3}$ e $V_{+U4}$ sono entrambi delle tensioni d'ingresso a degli \textit{inseguitori}
di tensione con rispettive uscite $V_A$ e $V_B$, infine dovr\'a valere (scritto in maniera elegante
ma forse meno leggibile)
\begin{equation*}
  V_A = -\max\{V_S, 0\} \qquad V_B = -\min\{V_S, 0\} \qquad V_A + V_B = -V_S
\end{equation*}

\paragraph{Richiesta (c)}
Determinare la funzione di trasferimento del secondo stadio del circuito (con ingressi VA e
VB ed uscita VOUT, verificando che si tratti di un amplificatore differenziale a guadagno
complesso).

\paragraph{Richiesta (d)}
Calcolare il modulo del guadagno di modo differenziale atteso nel limite di continua ($f\to 0$)
e ad 1kHz.

\paragraph{Richiesta (e)}
Calcolare il valore atteso per il rapporto tra il livello di continua di VOUT e l’ampiezza di VS
e lo si confronti con la pendenza della retta di best-fit ottenuta al punto 1d.

\paragraph{Richiesta (f)}
Si dia una stima del valore atteso percentuale dell’ampiezza della prima
armonica di VOUT, ovvero del rapporto tra la sua escursione ed il suo valor medio, e lo si
metta in relazione con le misura di ripple di cui al punto 1b.


% \paragraph{Richiesta (a)}
Sulla base di semplici argomenti di natura qualitativa (ovvero discutendone la risposta ai
limiti del dominio delle frequenze reali senza effettuare il calcolo della funzione di
trasferimento, come richiesto invece al punto successivo), mostrare che sia il blocco 1
che l’intero circuito svolgano effettivamente la funzione di filtro passa-basso e calcolare
per ciascuno il guadagno nel limite di continua.

\paragraph{Richiesta (b)}
Derivare analiticamente le funzioni di trasferimento $A_1(s)=V_1(s)/V_S(s)$,
$A_2(s)=V_2(s)/V_1(s)$ nonché la funzione di trasferimento complessiva $A(s)$,
individuando per ciascuna poli e zeri (si assuma che gli operazionali siano
ideali e per semplicit\'a che R1=R2=R3=R5=R=10kW, R2=R/10=1kW e C1=C2=C=100nF),
e verificare che gli andamenti osservati al punto 1a siano in accordo.

\paragraph{Risposta (b)}
Un'amplificatore invertente con opamp e impedenze ha una soluzione ben nota
\begin{equation*}
  H_1(s) = -\frac{R_2 || \frac{1}{sC_1}}{R_1} = -\frac{1}{1 + sRC}
\end{equation*}
\begin{equation*}
  H_2(s) = -\frac{R_5}{\left(R_3 || \frac{1}{sC_2}\right) + R_4} = -10\frac{1 + sRC}{11 + sRC}
\end{equation*}
\begin{equation}
  \label{eq:tfunc-compensazione-polo-zero}
  \colorbox{blu}{\boxed{H(s) = \frac{10}{11 + sRC}}}
\end{equation}

\paragraph{Richiesta (c)}
Dalle espressioni ottenute dedurre i valori attesi per i parametri dei filtri realizzati
rispettivamente dal blocco 1 e dall’intero circuito (frequenze di taglio, guadagno
massimo) e confrontarli con le misure di cui al punto 1b.

\paragraph{Richiesta (d)}
Determinare l’andamento temporale atteso per i segnali in V1 e V2, verificando che
descrivano con sufficiente accuratezza le forme d’onda osservate al punto 1c.

\paragraph{Richiesta (e)}
Calcolare i valori attesi per il tempo di discesa di V1 e di salita di V2 e confrontarli con
quanto misurato al punto 1d.

\paragraph{Richiesta (f)}
Illustrare qualitativamente l’effetto, sia nel dominio temporale che delle frequenze,
dell’applicazione del blocco 2 al segnale in V1.




% Si assuma l'operazionale ideale con VOUT in regime lineare. Si considerino i valori delle
resistenze pari a quelli nominali (in particolare R1=R2). Si assuma inoltre per semplicit\'a
Vg=0.

\paragraph{Richiesta (a)}
Si discuta lo stato di polarizzazione del diodo in corrispondenza di ciascuna semi-onda di VS
e si commenti l'andamento osservato in V+.

\paragraph{Richiesta (b)}
Per ciascuno dei due stati si determini la relazione tra VOUT e VS e si giustifichi l'andamento
osservato per VOUT al punto 1a. Di che tipo di circuito si tratta?

\paragraph{Richiesta (c)}
Si discuta il tipo di filtro operato da R4-C1 e le caratteristiche del segnale in VM.

\paragraph{Richiesta (d)}
Si calcoli il valore atteso per il rapporto tra il livello di continua di VM e l'ampiezza di VS e
lo si confronti con la pendenza della retta di best-fit ottenuta al punto 1d.

\paragraph{Richiesta (e)}
Si dia una stima del valore atteso percentuale del ripple di VM, ovvero del
rapporto tra la sua escursione ed il suo valor medio.


% \section*{Instrumentation Amplifier}
Si assumano gli operazionali ideali e le loro uscite
lineari. Si considerino inoltre i valori delle resistenze
pari a quelli nominali (e quindi R1=…=R4=R=10 kW).

\paragraph{Richiesta (a)}
Nell'ipotesi che il ramo di RG sia aperto, si calcoli il guadagno di modo differenziale
dell'amplificatore e se ne verifichi l'accordo con quanto ottenuto al punto 1b (suggerimento:
si consiglia di sfruttare la validità dei teoremi per le reti lineari e/o ricondurre il circuito ad
uno schema a blocchi di funzioni già note).

\paragraph{Richiesta (b)}
Si mostri che per ogni RG si abbia VOUT=0 se V1=V2, ovvero che il valore atteso per il
guadagno di modo comune Ac sia nullo.

\paragraph{Richiesta (c)}
Si calcoli più in generale l'espressione del guadagno di modo differenziale per valori finiti
di RG esplicitandone la dipendenza dal rapporto adimensionale $k=R/RG$.

\paragraph{Richiesta (d)}
Si calcoli quindi il valore atteso di Ad per RG=10, 5kW ($\implies k=1,2$) e lo si confronti con
quanto ottenuto al punto 2b.


% \input{src/ex_transconduttanza.tex}

% \input{src/ex_pompa_howland.tex}

\section*{Metodi}

\begin{equation}
  \begin{gathered}
    p(s) = \sum_{i=0}^n a_is^i \implies
    \begin{cases}
      \mathcal{R}_{\omega}[p] \equiv \sum_{i=0}^{\lfloor n/2\rfloor} (-1)^ia_{2i}\omega^{2i} \\
      \mathcal{I}_{\omega}[p] \equiv \sum_{i=1}^{\lfloor n/2\rfloor} (-1)^{i+1}a_{2i+1}\omega^{2i+1}
    \end{cases}
  \end{gathered}
\end{equation}

\paragraph{Forma canonica}
\begin{equation}
  \label{eq:forma-canonica}
  \begin{gathered}
    H(s) = \frac{\sum_i a_is^i}{\sum_i b_is^i} = \frac{N(s)}{D(s)} = \frac{\mathcal{R}_{\omega}^N + j\mathcal{I}_{\omega}^N}{\mathcal{R}_{\omega}^D + j\mathcal{I}_{\omega}^D} \\
    = \frac{(\mathcal{R}_{\omega}^N + j\mathcal{I}_{\omega}^N)(\mathcal{R}_{\omega}^D - j\mathcal{I}_{\omega}^D)}{(\mathcal{R}_{\omega}^D)^2 + (\mathcal{I}_{\omega}^D)^2} \\
    = \left(\frac{(\mathcal{R}_{\omega}^N\mathcal{R}_{\omega}^D + \mathcal{I}_{\omega}^N\mathcal{I}_{\omega}^D)}{(\mathcal{R}_{\omega}^D)^2 + (\mathcal{I}_{\omega}^D)^2}\right) + j\left(\frac{(\mathcal{I}_{\omega}^N\mathcal{R}_{\omega}^D - \mathcal{R}_{\omega}^N\mathcal{I}_{\omega}^D)}{(\mathcal{R}_{\omega}^D)^2 + (\mathcal{I}_{\omega}^D)^2}\right)
  \end{gathered}
\end{equation}

\paragraph{Guadagno}
\begin{equation}
  \label{eq:da-tfunc-a-gain}
  \begin{gathered}
    G(\omega)^2 \equiv |H(j\omega)|^2 = H(j\omega)\overline{H(j\omega)} = \\
    \frac{(\mathcal{R}_{\omega}^N\mathcal{R}_{\omega}^D + \mathcal{I}_{\omega}^N(\mathcal{I}_{\omega}^D))^2 + (\mathcal{I}_{\omega}^N\mathcal{R}_{\omega}^D - \mathcal{R}_{\omega}^N(\mathcal{I}_{\omega}^D))^2}{((\mathcal{R}_{\omega}^D)^2 + (\mathcal{I}_{\omega}^D)^2)^2} =
    \frac{(\mathcal{R}_{\omega}^N)^2(\mathcal{R}_{\omega}^D)^2 + (\mathcal{I}_{\omega}^N)^2(\mathcal{I}_{\omega}^D)^2 + (\mathcal{I}_{\omega}^N)^2(\mathcal{R}_{\omega}^D)^2 + (\mathcal{R}_{\omega}^N)^2(\mathcal{I}_{\omega}^D)^2}{((\mathcal{R}_{\omega}^D)^2 + (\mathcal{I}_{\omega}^D)^2)^2} = \\
    \frac{((\mathcal{R}_{\omega}^N)^2 + (\mathcal{I}_{\omega}^N)^2)((\mathcal{R}_{\omega}^D)^2 + (\mathcal{I}_{\omega}^D)^2)}{((\mathcal{R}_{\omega}^D)^2 + (\mathcal{I}_{\omega}^D)^2)^2} = \frac{(\mathcal{R}_{\omega}^N)^2 + (\mathcal{I}_{\omega}^N)^2}{(\mathcal{R}_{\omega}^D)^2 + (\mathcal{I}_{\omega}^D)^2}
  \end{gathered}
\end{equation}

\paragraph{Fase}
\begin{itemize}
\item Semipiano reale positivo $\mathcal{R}_{\omega}[H(s)] > 0$
  \begin{equation}
    \label{eq:da-tfunc-a-fase-reale-positivo}
    \begin{gathered}
      \mathcal{R}_{\omega}^N\mathcal{R}_{\omega}^D + \mathcal{I}_{\omega}^N\mathcal{I}_{\omega}^D > 0 \\ \implies
      \phi(\omega) = \arctan\left(\frac{\mathcal{I}_{\omega}^N\mathcal{R}_{\omega}^D - \mathcal{R}_{\omega}^N\mathcal{I}_{\omega}^D}{\mathcal{R}_{\omega}^N\mathcal{R}_{\omega}^D + \mathcal{I}_{\omega}^N\mathcal{I}_{\omega}^D}\right)
    \end{gathered}
  \end{equation}
\item Semipiano immaginario positivo $\mathcal{I}_{\omega}[H(s)] > 0$
  \begin{equation}
    \label{eq:da-tfunc-a-fase-immaginario-positivo}
    \begin{gathered}
      \mathcal{I}_{\omega}^N\mathcal{R}_{\omega}^D - \mathcal{R}_{\omega}^N\mathcal{I}_{\omega}^D > 0 \\ \implies
      \phi(\omega) = \frac{\pi}{2} + \arctan\left(-\frac{\mathcal{R}_{\omega}^N\mathcal{R}_{\omega}^D + \mathcal{I}_{\omega}^N\mathcal{I}_{\omega}^D}{\mathcal{I}_{\omega}^N\mathcal{R}_{\omega}^D - \mathcal{R}_{\omega}^N\mathcal{I}_{\omega}^D}\right)
    \end{gathered}
  \end{equation}

\item Semipiano reale negativo $\mathcal{R}_{\omega}[H(s)] < 0$
  \begin{equation}
    \label{eq:da-tfunc-a-fase-reale-negativo}
    \begin{gathered}
      \mathcal{R}_{\omega}^N\mathcal{R}_{\omega}^D + \mathcal{I}_{\omega}^N\mathcal{I}_{\omega}^D < 0 \\ \implies
      \phi(\omega) = \pm\pi + \arctan\left(\frac{\mathcal{I}_{\omega}^N\mathcal{R}_{\omega}^D - \mathcal{R}_{\omega}^N\mathcal{I}_{\omega}^D}{\mathcal{R}_{\omega}^N\mathcal{R}_{\omega}^D + \mathcal{I}_{\omega}^N\mathcal{I}_{\omega}^D}\right)
    \end{gathered}
  \end{equation}
\item Semipiano immaginario negativo $\mathcal{I}_{\omega}[H(s)] < 0$
  \begin{equation}
    \label{eq:da-tfunc-a-fase-immaginario-negativo}
    \begin{gathered}
      \mathcal{I}_{\omega}^N\mathcal{R}_{\omega}^D - \mathcal{R}_{\omega}^N\mathcal{I}_{\omega}^D < 0 \\ \implies
      \phi(\omega) = -\frac{\pi}{2} + \arctan\left(-\frac{\mathcal{R}_{\omega}^N\mathcal{R}_{\omega}^D + \mathcal{I}_{\omega}^N\mathcal{I}_{\omega}^D}{\mathcal{I}_{\omega}^N\mathcal{R}_{\omega}^D - \mathcal{R}_{\omega}^N\mathcal{I}_{\omega}^D}\right)
    \end{gathered}
  \end{equation}
\end{itemize}


\pagebreak

\end{document}
